\documentclass[11pt,a4paper]{article}

% ==================== TEMEL PAKETLER (pdflatex) ====================
\usepackage[T1]{fontenc}
\usepackage[utf8]{inputenc}
\usepackage{lmodern}
\usepackage[a4paper,margin=2cm,top=2.4cm,bottom=2.4cm]{geometry}
\usepackage{xcolor}
\usepackage{colortbl}
\usepackage{array}
\usepackage{booktabs}
\usepackage{longtable}
\usepackage{tabularx}
\usepackage{listings}
\usepackage{fancyhdr}
\usepackage{amsmath}
\usepackage{amssymb}
\usepackage{hyperref}

% ==================== RENK PALETİ ====================
\definecolor{anaMavi}{HTML}{1B3A5B}
\definecolor{vurguMavi}{HTML}{2E6DA4}
\definecolor{acikMavi}{HTML}{EAF2F9}
\definecolor{turuncu}{HTML}{D9701E}
\definecolor{yesil}{HTML}{2E7D32}
\definecolor{acikYesil}{HTML}{E8F5E9}
\definecolor{kirmizi}{HTML}{C0392B}
\definecolor{mor}{HTML}{6C3483}
\definecolor{acikMor}{HTML}{F3E9F7}
\definecolor{acikGri}{HTML}{F4F5F7}
\definecolor{ortaGri}{HTML}{6B7280}
\definecolor{koyuGri}{HTML}{2D2D2D}
\definecolor{kodArka}{HTML}{FBFBFA}

% ==================== HYPERREF ====================
\hypersetup{
  colorlinks=true,
  linkcolor={[HTML]{2E6DA4}},
  urlcolor={[HTML]{2E6DA4}},
  pdftitle={Python ile Veri Bilimi Rehberi},
  pdfauthor={Veri Bilimi Rehberi},
  bookmarksnumbered=true
}

% ==================== BAŞLIK STİLLERİ (titlesec olmadan) ====================
\setcounter{secnumdepth}{0}
\newcounter{secc}
\newcounter{subsecc}[secc]
\newcounter{subsubsecc}[subsecc]
\newcommand{\gsec}[1]{\par\phantomsection\stepcounter{secc}%
  \setcounter{subsecc}{0}%
  \vspace{18pt}{\color{anaMavi}\Large\bfseries\thesecc.\hspace{0.5em}#1}\par
  \vspace{2pt}{\color{turuncu}\rule{\linewidth}{1.6pt}}\par\vspace{8pt}%
  \addcontentsline{toc}{section}{\protect\numberline{\thesecc}#1}\markboth{#1}{#1}}
\newcommand{\gsub}[1]{\par\phantomsection\stepcounter{subsecc}%
  \setcounter{subsubsecc}{0}%
  \vspace{11pt}{\color{vurguMavi}\large\bfseries\thesecc.\thesubsecc\hspace{0.5em}#1}\par\vspace{4pt}%
  \addcontentsline{toc}{subsection}{\protect\numberline{\thesecc.\thesubsecc}#1}}
\newcommand{\gsubsub}[1]{\par\phantomsection\stepcounter{subsubsecc}%
  \vspace{7pt}{\color{koyuGri}\normalsize\bfseries #1}\par\vspace{2pt}}

% Kısım (Part) başlığı
\newcommand{\gpart}[2]{\clearpage
  \begin{center}\vspace*{1.5cm}
  {\color{turuncu}\rule{0.7\linewidth}{1.5pt}}\\[10pt]
  {\color{ortaGri}\Large KISIM #1}\\[8pt]
  {\color{anaMavi}\Huge\bfseries #2}\\[10pt]
  {\color{turuncu}\rule{0.7\linewidth}{1.5pt}}
  \end{center}\vspace{1cm}}

% ==================== BAŞLIK / ALTBİLGİ ====================
\setlength{\headheight}{14pt}
\pagestyle{fancy}
\fancyhf{}
\renewcommand{\headrulewidth}{0.4pt}
\renewcommand{\footrulewidth}{0.4pt}
\fancyhead[L]{\small\color{ortaGri}Python ile Veri Bilimi}
\fancyhead[R]{\small\color{ortaGri}\leftmark}
\fancyfoot[C]{\small\color{ortaGri}\thepage}

% ==================== KOD LİSTELEME (Python) ====================
\lstdefinestyle{pystyle}{
  language=Python,
  basicstyle=\ttfamily\small\color{koyuGri},
  keywordstyle=\color{vurguMavi}\bfseries,
  commentstyle=\color{yesil}\itshape,
  stringstyle=\color{turuncu},
  numberstyle=\tiny\color{ortaGri},
  numbers=left,
  numbersep=8pt,
  backgroundcolor=\color{kodArka},
  frame=single,
  framesep=6pt,
  rulecolor=\color{ortaGri!40},
  breaklines=true,
  showstringspaces=false,
  tabsize=4,
  columns=fullflexible,
  keepspaces=true,
  extendedchars=true,
  literate={ı}{{\i}}1{İ}{{\.I}}1{ş}{{\c s}}1{Ş}{{\c S}}1{ğ}{{\u g}}1{Ğ}{{\u G}}1{ç}{{\c c}}1{Ç}{{\c C}}1{ö}{{\"o}}1{Ö}{{\"O}}1{ü}{{\"u}}1{Ü}{{\"U}}1,
  morekeywords={np,pd,plt,sns,as,with,yield,lambda,None,True,False,
    self,assert,async,await,nonlocal,DataFrame,Series,array,ndarray,
    fit,transform,predict,fit_transform}
}
\lstset{style=pystyle}

% Kabuk/terminal stili
\lstdefinestyle{shstyle}{
  basicstyle=\ttfamily\small\color{koyuGri},
  backgroundcolor=\color{koyuGri!6},
  frame=single,framesep=6pt,rulecolor=\color{ortaGri!40},
  breaklines=true,showstringspaces=false,tabsize=4,columns=fullflexible,
  keepspaces=true,
  literate={ı}{{\i}}1{İ}{{\.I}}1{ş}{{\c s}}1{Ş}{{\c S}}1{ğ}{{\u g}}1{Ğ}{{\u G}}1{ç}{{\c c}}1{Ç}{{\c C}}1{ö}{{\"o}}1{Ö}{{\"O}}1{ü}{{\"u}}1{Ü}{{\"U}}1,
}

% ==================== ÖZEL KUTULAR (tcolorbox olmadan) ====================
\newcommand{\callout}[4]{\par\smallskip\noindent\begingroup
  \setlength{\fboxsep}{8pt}\setlength{\fboxrule}{0.8pt}%
  \fcolorbox{#1}{#2}{\parbox{\dimexpr\linewidth-2\fboxsep-2\fboxrule\relax}{%
    {\bfseries\color{#1}#3}\par\smallskip #4}}\endgroup\par\smallskip}
\newcommand{\bilgi}[1]{\callout{vurguMavi}{acikMavi}{Bilgi}{#1}}
\newcommand{\ipucu}[1]{\callout{yesil}{acikYesil}{İpucu}{#1}}
\newcommand{\uyari}[1]{\callout{kirmizi}{kirmizi!7}{Dikkat}{#1}}
\newcommand{\ornek}[2][Örnek]{\callout{ortaGri!90!black}{acikGri}{#1}{#2}}
\newcommand{\notk}[2][Not]{\callout{mor}{acikMor}{#1}{#2}}

% Yardımcı komutlar
\newcommand{\code}[1]{\texttt{\color{koyuGri}#1}}
\renewcommand{\arraystretch}{1.35}

% ==================== BAŞLANGIÇ ====================
\begin{document}

% ---------- KAPAK ----------
\begin{titlepage}
\centering
\vspace*{1.4cm}
{\color{turuncu}\rule{\textwidth}{2pt}}\\[8pt]
{\color{anaMavi}\Huge\bfseries Python ile Veri Bilimi}\\[10pt]
{\color{vurguMavi}\LARGE Teoriden Uygulamaya Kapsamlı Rehber}\\[8pt]
{\color{ortaGri}\large NumPy \textbullet\ Pandas \textbullet\ Matplotlib \textbullet\ scikit-learn}\\[6pt]
{\color{turuncu}\rule{\textwidth}{2pt}}\\[1.4cm]

\begin{center}
\fcolorbox{vurguMavi}{acikMavi}{\parbox{0.86\textwidth}{\centering\large
Bu belge, veri bilimini Python ekosistemiyle --- \textbf{teorisi ve çalışan
kod örnekleriyle} --- baştan sona anlatır. Sayısal hesaplamadan (NumPy) veri
işlemeye (Pandas), görselleştirmeden (Matplotlib/Seaborn) istatistik ve makine
öğrenmesine (scikit-learn) kadar tam bir veri bilimi hattını kapsar.\vspace{4pt}}}
\end{center}

\vspace{0.8cm}
\begin{center}
\fcolorbox{ortaGri!50}{acikGri}{\parbox{0.86\textwidth}{\small
\textbf{Kapsam --- Kısım 1 (Temeller \& NumPy):} veri bilimi iş akışı
\textbullet\ Python ekosistemi \textbullet\ \code{ndarray} \textbullet\ dtype
\textbullet\ yayınlama (broadcasting) \textbullet\ vektörleştirme \textbullet\
indeksleme \textbullet\ doğrusal cebir \& rastgelelik.\vspace{4pt}

\textbf{Kapsam --- Kısım 2 (Pandas):} \code{Series} \& \code{DataFrame}
\textbullet\ veri okuma/yazma \textbullet\ seçme \& filtreleme \textbullet\
veri temizleme \textbullet\ \code{groupby} \& pivot \textbullet\ birleştirme
\textbullet\ zaman serileri.\vspace{4pt}

\textbf{Kapsam --- Kısım 3 (Görselleştirme):} Matplotlib temelleri \textbullet\
alt grafikler \& stil \textbullet\ Seaborn ile istatistiksel grafikler.\vspace{4pt}

\textbf{Kapsam --- Kısım 4 (İstatistik):} betimsel istatistik \textbullet\
olasılık dağılımları \textbullet\ korelasyon \textbullet\ hipotez testleri.\vspace{4pt}

\textbf{Kapsam --- Kısım 5 (Makine Öğrenmesi):} veri hazırlama \textbullet\
regresyon \textbullet\ sınıflandırma \textbullet\ kümeleme \textbullet\
model değerlendirme \textbullet\ pipeline'lar.\vspace{4pt}}}
\end{center}

\vfill
{\color{ortaGri}\large Hazırlanma Tarihi: 12 Temmuz 2026}\\[4pt]
{\color{ortaGri} Python 3 \textbullet\ NumPy \textbullet\ Pandas \textbullet\ scikit-learn}
\vspace*{0.6cm}
\end{titlepage}

% ---------- İÇİNDEKİLER ----------
\tableofcontents
\thispagestyle{empty}
\clearpage

% #####################################################################
\gpart{1}{Temeller ve NumPy}
% #####################################################################

% =====================================================================
\gsec{Veri Bilimine Giriş}
% =====================================================================

Veri bilimi, ham veriden anlamlı bilgi ve öngörü çıkarma disiplinidir;
istatistik, programlama ve alan bilgisinin kesişiminde durur. Bir veri bilimi
projesi nadiren ``model eğitmek''ten ibarettir; zamanın büyük kısmı veriyi
toplamak, temizlemek ve anlamakla geçer.

\gsub{Veri Bilimi İş Akışı}

Tipik bir proje şu aşamalardan geçer. Bu rehber, her aşama için gereken
araçları sırayla öğretir.

\begin{center}
\begin{tabularx}{\textwidth}{@{}>{\bfseries\color{anaMavi}}l X c@{}}
\toprule
\rowcolor{acikMavi}\color{anaMavi}\textbf{Aşama} & \color{anaMavi}\textbf{Ne yapılır?} & \color{anaMavi}\textbf{Araç}\\
\midrule
1. Toplama & Veriyi kaynaklardan al (CSV, SQL, API) & Pandas\\
\rowcolor{acikGri} 2. Temizleme & Eksik/hatalı/tekrar veriyi düzelt & Pandas\\
3. Keşif (EDA) & İstatistik ve grafiklerle veriyi anla & Pandas, Matplotlib\\
\rowcolor{acikGri} 4. Modelleme & İstatistik/ML modeli kur & scikit-learn\\
5. Değerlendirme & Modeli ölç, yorumla & scikit-learn\\
\rowcolor{acikGri} 6. İletişim & Sonuçları görselleştir, sun & Matplotlib, Seaborn\\
\bottomrule
\end{tabularx}
\end{center}

\gsub{Python Veri Bilimi Ekosistemi}

\begin{center}
\begin{tabularx}{\textwidth}{@{}>{\ttfamily\bfseries\color{anaMavi}}l X@{}}
\toprule
\rowcolor{acikMavi}\color{anaMavi}\textbf{Kütüphane} & \color{anaMavi}\textbf{Rol}\\
\midrule
NumPy & Sayısal hesaplama; çok boyutlu diziler, temel; her şey buna dayanır.\\
\rowcolor{acikGri} Pandas & Etiketli tablo verisi (DataFrame); okuma, temizleme, dönüştürme.\\
Matplotlib & Temel çizim kütüphanesi; her tür grafik.\\
\rowcolor{acikGri} Seaborn & Matplotlib üzerine istatistiksel, güzel grafikler.\\
scikit-learn & Makine öğrenmesi; regresyon, sınıflandırma, kümeleme.\\
\rowcolor{acikGri} SciPy & İleri bilimsel hesaplama; istatistik, optimizasyon.\\
\bottomrule
\end{tabularx}
\end{center}

\gsub{Kurulum ve Ortam}

\begin{lstlisting}[style=shstyle]
# Sanal ortam olustur (izole, temiz kurulum)
$ python -m venv datascience
$ source datascience/bin/activate       # Linux/Mac
# datascience\Scripts\activate          # Windows

# Temel kutuphaneleri kur
$ pip install numpy pandas matplotlib seaborn scikit-learn jupyter

# Jupyter Notebook baslat (interaktif calisma icin ideal)
$ jupyter notebook
\end{lstlisting}

\bilgi{Veri bilimi çalışmaları için \textbf{Jupyter Notebook} (veya JupyterLab)
neredeyse standarttır: kodu hücre hücre çalıştırır, grafikleri satır içinde
gösterir ve keşifsel analizi çok kolaylaştırır. Bu rehberdeki tüm örnekler bir
Jupyter hücresinde veya bir \code{.py} dosyasında çalıştırılabilir.}

% =====================================================================
\gsec{NumPy: Sayısal Hesaplamanın Temeli}
% =====================================================================

NumPy (Numerical Python), tüm veri bilimi ekosisteminin altında yatan
kütüphanedir. Kalbi \code{ndarray} (N-boyutlu dizi) nesnesidir: aynı tipteki
verileri bellekte ardışık tutan, üzerinde \emph{vektörleştirilmiş} (döngüsüz)
işlemler yapılabilen hızlı bir dizi.

\gsub{Neden NumPy? Liste vs Dizi}

Python listeleri esnektir ama yavaştır: her eleman ayrı bir nesnedir ve döngüler
Python yorumlayıcısında çalışır. NumPy dizileri veriyi bitişik bellekte tutar
ve işlemleri derlenmiş C koduyla yapar --- genelde 10--100 kat daha hızlı.

\begin{lstlisting}
import numpy as np

# Python listesi ile: yavas, dongu gerekir
lst = [1, 2, 3, 4]
squares = [x**2 for x in lst]            # [1, 4, 9, 16]

# NumPy ile: hizli, vektorlestirilmis (dongu YOK)
arr = np.array([1, 2, 3, 4])
squares = arr ** 2                        # array([1, 4, 9, 16])
\end{lstlisting}

\gsub{Dizi Oluşturma}

\begin{lstlisting}
import numpy as np

a = np.array([1, 2, 3])                   # listeden
b = np.zeros((2, 3))                      # 2x3 sifir matrisi
c = np.ones((3,))                         # [1., 1., 1.]
d = np.full((2, 2), 7)                    # 7 ile dolu 2x2
e = np.eye(3)                             # 3x3 birim matris
f = np.arange(0, 10, 2)                   # [0 2 4 6 8]
g = np.linspace(0, 1, 5)                  # [0. 0.25 0.5 0.75 1.]
h = np.random.rand(2, 3)                  # 2x3 rastgele [0,1)
\end{lstlisting}

\gsub{Dizi Özellikleri: shape, dtype, boyut}

Bir dizinin \emph{şekli} (shape) boyutlarını, \emph{dtype}'ı ise sakladığı
veri tipini belirtir. Bunları anlamak NumPy'da hataların çoğunu önler.

\begin{lstlisting}
a = np.array([[1, 2, 3],
              [4, 5, 6]])
a.shape        # (2, 3) -> 2 satir, 3 sutun
a.ndim         # 2      -> boyut sayisi
a.size         # 6      -> toplam eleman
a.dtype        # dtype('int64')

# dtype donusumu (bellek/hassasiyet kontrolu icin onemli)
b = a.astype(np.float32)
c = np.array([1, 2, 3], dtype=np.float64)
\end{lstlisting}

\gsub{Yeniden Şekillendirme (Reshape)}

\begin{lstlisting}
a = np.arange(12)                         # [0 1 2 ... 11]
b = a.reshape(3, 4)                       # 3x4 matrise cevir
c = a.reshape(2, 2, 3)                    # 3 boyutlu
d = b.reshape(-1)                         # duzlestir -> 1B (12,)
e = b.T                                   # transpoz (4x3)
f = b.reshape(4, -1)                      # -1: boyutu otomatik hesapla (4x3)
\end{lstlisting}

\uyari{\code{reshape} verinin bir \emph{görünümünü} (view) döndürür, kopya
değil; yani yeni diziyi değiştirmek çoğu zaman orijinali de değiştirir. Bağımsız
bir kopya istiyorsan \code{a.copy()} kullan. Bu view/copy ayrımı, NumPy'da
şaşırtıcı hataların en yaygın kaynağıdır.}

% =====================================================================
\gsec{NumPy: İndeksleme, Yayınlama ve İşlemler}
% =====================================================================

\gsub{İndeksleme ve Dilimleme}

NumPy, Python listelerinden daha güçlü indeksleme sunar: çok boyutlu dilimleme,
boole maskeleri ve süslü (fancy) indeksleme.

\begin{lstlisting}
a = np.array([[1, 2, 3, 4],
              [5, 6, 7, 8],
              [9, 10, 11, 12]])

a[0, 2]        # 3      -> satir 0, sutun 2
a[1]           # [5 6 7 8]  -> tum 1. satir
a[:, 1]        # [2 6 10]   -> tum satirlarin 1. sutunu
a[0:2, 1:3]    # [[2 3] [6 7]]  -> alt matris
a[-1]          # [9 10 11 12]   -> son satir

# Boole maskesi (kosula uyan elemanlar) -- cok guclu!
a[a > 6]       # [7 8 9 10 11 12]
a[a % 2 == 0]  # cift elemanlar

# Susulu (fancy) indeksleme
a[[0, 2]]      # 0. ve 2. satirlar
\end{lstlisting}

\gsub{Yayınlama (Broadcasting)}

Yayınlama, NumPy'ın \emph{farklı şekilli} dizileri birlikte işlemesini sağlayan
kuraldır. Küçük dizi, büyük diziye ``yayılarak'' uyumlu hale getirilir ---
bellekte kopya oluşturmadan. Bu, veri biliminde standardizasyon gibi işlemlerin
temelidir.

\begin{lstlisting}
a = np.array([[1, 2, 3],
              [4, 5, 6]])                 # (2, 3)

a + 10          # her elemana 10 ekle (skaler yayilir)
a * 2           # her elemani 2 ile carp

row = np.array([10, 20, 30])              # (3,)
a + row         # her SATIRA ekle -> [[11 22 33] [14 25 36]]

col = np.array([[100], [200]])            # (2, 1)
a + col         # her SUTUNA ekle -> [[101 102 103] [204 205 206]]

# Pratik: sutunlari standartlastir (z-skoru)
mean = a.mean(axis=0)                     # her sutunun ortalamasi
std = a.std(axis=0)
z = (a - mean) / std                      # yayinlama ile
\end{lstlisting}

\notk[Yayınlama kuralı]{İki dizi \emph{sondan başa} boyut boyut karşılaştırılır.
İki boyut uyumludur eğer: (1) eşitlerse, ya da (2) biri 1 ise. Örneğin (2,3)
ile (3,) uyumludur çünkü (3,), (1,3) gibi ele alınır ve satır boyunca yayılır.}

\gsub{Eksen (axis) Boyunca İşlemler}

Toplama, ortalama gibi indirgeme (reduction) işlemleri bir \emph{eksen} boyunca
yapılabilir. \code{axis=0} sütunlar boyunca (satırları çökertir), \code{axis=1}
satırlar boyunca çalışır. Bu kavram Pandas'ta da aynen geçerlidir.

\begin{lstlisting}
a = np.array([[1, 2, 3],
              [4, 5, 6]])

a.sum()             # 21   -> tum elemanlar
a.sum(axis=0)       # [5 7 9]   -> her SUTUN toplami
a.sum(axis=1)       # [6 15]    -> her SATIR toplami
a.mean(axis=0)      # [2.5 3.5 4.5]
a.max(axis=1)       # [3 6]
a.argmax()          # 5   -> en buyugun (duzlestirilmis) indeksi
np.cumsum(a)        # kumulatif toplam
\end{lstlisting}

\gsub{Doğrusal Cebir}

NumPy, makine öğrenmesinin altında yatan matris işlemlerini sağlar.

\begin{lstlisting}
A = np.array([[1, 2],
              [3, 4]])
B = np.array([[5, 6],
              [7, 8]])

A @ B                    # matris carpimi (dot product)
A * B                    # ELEMAN BAZLI carpim (dikkat: farkli!)
np.dot(A, B)             # A @ B ile ayni
A.T                      # transpoz
np.linalg.inv(A)         # ters matris
np.linalg.det(A)         # determinant
eigenvalues, eigenvectors = np.linalg.eig(A)  # ozdegerler/ozvektorler

# Dogrusal denklem sistemi cozme: Ax = b
b = np.array([1, 2])
x = np.linalg.solve(A, b)
\end{lstlisting}

\gsub{Rastgelelik ve Tekrarlanabilirlik}

\begin{lstlisting}
rng = np.random.default_rng(seed=42)      # tekrarlanabilir (modern API)

rng.random(3)                # [0,1) duzgun dagilim
rng.integers(1, 7, size=10)  # 1-6 arasi 10 zar atisi
rng.normal(0, 1, size=5)     # standart normal dagilim
rng.choice([1, 2, 3], size=4)# ornekleme
shuffled = rng.permutation(np.arange(10))# karistirma
\end{lstlisting}

\ipucu{\textbf{Her zaman bir seed (tohum) belirle} (\code{default\_rng(42)}
veya \code{np.random.seed(42)}). Böylece analizlerin ve model sonuçların
tekrarlanabilir olur --- kod her çalıştığında aynı ``rastgele'' değerleri
üretir, bu da hata ayıklamayı ve sonuçları paylaşmayı kolaylaştırır.}

% =====================================================================
\gsec{NumPy Derinlemesine: Veri Tipleri ve Bellek}
% =====================================================================

Yukarıdaki temeller günlük işin \%80'ini görür. Ama NumPy'ı gerçekten
\emph{anlamak} --- neden hızlı olduğunu, neden bazen şaşırtıcı sonuçlar
verdiğini --- birkaç düşük seviye kavramı bilmeyi gerektirir. Bu bölüm o
kavramları bol örnekle açar.

\gsub{Veri Tipleri (dtype) Derinlemesine}

Bir NumPy dizisindeki \emph{tüm} elemanlar aynı tiptedir ve her tip bellekte
sabit sayıda bayt kaplar. Doğru dtype seçmek hem bellek hem hız kazandırır:
100 milyonluk bir diziyi \code{int64} yerine \code{int8} tutmak belleği 8 kat
azaltır.

\begin{lstlisting}
import numpy as np

# Tam sayilar: int8/16/32/64 (isaretli), uint8/... (isaretsiz)
a = np.array([1, 2, 3], dtype=np.int8)    # her eleman 1 bayt
a.itemsize                                 # 1  -> eleman basina bayt
a.nbytes                                   # 3  -> toplam bayt

# Ondalikli: float16/32/64
b = np.array([1.5, 2.5], dtype=np.float32)

# Bool ve metin
c = np.array([True, False, True])          # dtype('bool')
d = np.array(['al', 'sat'])                # dtype('<U3') sabit uzunluk!

# Karisik listeden olusturmada NumPy ortak tip secer (upcasting)
np.array([1, 2, 3.0]).dtype                # float64 (int -> float'a yukselir)
np.array([1, 'x']).dtype                   # '<U21' (hepsi metne dondu!)
\end{lstlisting}

\uyari{\textbf{Tam sayı taşması (overflow) sessizdir!} Sabit boyutlu tipler
sınırı aşınca döner (wrap-around) --- hata vermez. \code{int8} en fazla 127
tutar; \code{np.int8(127) + 1} sonucu \code{-128} olur. Büyük değerlerle
çalışırken \code{int64} veya \code{float64} kullan.}

\begin{lstlisting}
x = np.array([127], dtype=np.int8)
x + 1                    # array([-128])  <- SESSIZ TASMA!

# astype ile guvenli donusum
big = x.astype(np.int64) + 1   # array([128])  dogru

# float -> int donusumunde ondalik ATILIR (yuvarlanmaz)
np.array([1.9, 2.1, -1.9]).astype(int)     # array([ 1,  2, -1])
np.round([1.9, 2.1]).astype(int)           # array([2, 2])  once yuvarla
\end{lstlisting}

\gsub{Görünüm (View) mü, Kopya (Copy) mı?}

Bu, NumPy'daki en sık ``neden değişti?!'' hatasının kaynağıdır. \textbf{Dilimleme
(slicing) bir görünüm döndürür} --- aynı belleğe bakan yeni bir pencere.
Görünümü değiştirmek orijinali de değiştirir. \textbf{Süslü/boole indeksleme ise
kopya döndürür.}

\begin{lstlisting}
a = np.arange(10)                # [0 1 2 3 4 5 6 7 8 9]

# DILIM = GORUNUM (ayni bellek)
s = a[2:5]                       # [2 3 4]
s[0] = 999
a                                # [0 1 999 3 4 5 6 7 8 9]  <- a DEGISTI!

# .copy() ile bagimsiz kopya
s2 = a[2:5].copy()
s2[0] = -1                       # a etkilenmez

# Bir dizinin gorunum mu oldugunu kontrol et
s.base is a                      # True  -> s, a'nin gorunumu
s2.base is a                     # False -> bagimsiz kopya

# SUSLU/BOOLE indeksleme her zaman KOPYA doner
f = a[[1, 3, 5]]                 # kopya
m = a[a > 5]                     # kopya
f[0] = 0                         # a'yi ETKILEMEZ
\end{lstlisting}

\notk[Neden önemli?]{Fonksiyona dizi geçirip içinde dilimleyerek değiştirirsen,
çağıranın verisini farkında olmadan bozabilirsin. Kural: veriyi değiştirecek bir
fonksiyon yazıyorsan ya baştan \code{arr.copy()} al, ya da davranışı belgele.
Tersine, büyük veride gereksiz \code{.copy()} bellek israfıdır --- görünümler
bedavadır.}

\gsub{Bir Diziyi Yerinde Değiştirmek}

Görünümler, koşullu güncellemeyi çok verimli kılar: yeni dizi oluşturmadan.

\begin{lstlisting}
a = np.array([-3, 5, -1, 8, -7, 2])

# Boole maske ile YERINDE atama (negatifleri sifirla)
a[a < 0] = 0                     # [0 5 0 8 0 2]

# np.clip: degerleri bir araliga sikistir
np.clip(a, 0, 5)                 # [0 5 0 5 0 2]

# np.where: kosullu SECIM (if-else vektorlestirilmis)
b = np.array([10, -2, 7, -5])
np.where(b > 0, b, 0)           # pozitifse kendisi, degilse 0
np.where(b > 0, 'arti', 'eksi')# [arti eksi arti eksi]

# np.select: cok kosullu (elif zinciri gibi)
notlar = np.array([95, 82, 60, 45])
kosullar = [notlar >= 90, notlar >= 70, notlar >= 50]
sonuc = np.select(kosullar, ['A', 'B', 'C'], default='F')
# ['A' 'B' 'C' 'F']
\end{lstlisting}

% =====================================================================
\gsec{NumPy Derinlemesine: Vektörizasyon ve Performans}
% =====================================================================

\gsub{Evrensel Fonksiyonlar (ufunc) ve Vektörizasyon}

NumPy'ın gücü \emph{vektörizasyondan} gelir: bir işlemi tüm diziye tek seferde,
derlenmiş C döngüsüyle uygulamak. Python \code{for} döngüsünden kaçındığın her
yerde büyük hız kazanırsın. \code{np.sin}, \code{+}, \code{**} gibi eleman bazlı
çalışan fonksiyonlara \emph{ufunc} denir.

\begin{lstlisting}
x = np.linspace(0, 2*np.pi, 1_000_000)

# YAVAS: Python dongusu (kacin!)
# y = np.array([np.sin(v) for v in x])

# HIZLI: vektorlestirilmis (10-100x daha hizli)
y = np.sin(x)

# Eleman bazli ufunc'lar zincirlenebilir
z = np.exp(-x) * np.cos(x) + 1

# Matematiksel yardimcilar
np.sqrt([1, 4, 9])              # [1. 2. 3.]
np.log(np.e)                    # 1.0
np.abs([-2, 3, -5])             # [2 3 5]
np.maximum([1, 5], [4, 2])      # [4 5]  eleman bazli maks
\end{lstlisting}

\ornek[Örnek: Vektörizasyonun Gücü]{İki nokta kümesi arasındaki tüm mesafeleri
tek satırda, döngüsüz hesaplamak --- yayınlama + ufunc birleşimi:}

\begin{lstlisting}
noktalar = np.array([[0, 0], [3, 4], [6, 8]])   # (3, 2)

# Her ciftin oklit mesafesi (yayinlama ile)
fark = noktalar[:, None, :] - noktalar[None, :, :]   # (3, 3, 2)
mesafe = np.sqrt((fark ** 2).sum(axis=2))            # (3, 3)
# mesafe[i, j] = i. ve j. nokta arasi uzaklik
\end{lstlisting}

\bilgi{\code{np.vectorize} bir Python fonksiyonunu diziye uygular ama içeride
gerçek bir döngü çalıştırır --- \emph{hız kazandırmaz}, sadece kolaylık sağlar.
Gerçek hız için işlemi NumPy ufunc'larıyla ifade et.}

\gsub{Dizi Birleştirme ve Bölme}

Veriyi yan yana veya üst üste eklemek, ML'de özellik matrisleri kurmanın
temelidir.

\begin{lstlisting}
a = np.array([[1, 2], [3, 4]])
b = np.array([[5, 6], [7, 8]])

np.concatenate([a, b], axis=0)  # dikey (satir ekle) -> (4, 2)
np.concatenate([a, b], axis=1)  # yatay (sutun ekle) -> (2, 4)
np.vstack([a, b])               # dikey (v = vertical)
np.hstack([a, b])               # yatay (h = horizontal)

# stack: YENI bir eksen ekler
np.stack([a, b])                # (2, 2, 2) -> ustuste yeni katman

# Bir sutunu/satiri diziye ekleme icin newaxis
v = np.array([1, 2])            # (2,)
v[:, None]                      # (2, 1) sutun vektoru
v[None, :]                      # (1, 2) satir vektoru

# Bolme
c = np.arange(12).reshape(3, 4)
np.split(c, 2, axis=1)          # 4 sutunu 2 esit parcaya ([:,:2],[:,2:])
\end{lstlisting}

\gsub{Sıralama, Arama ve Faydalı İşlemler}

\begin{lstlisting}
a = np.array([3, 1, 4, 1, 5, 9, 2, 6])

np.sort(a)                      # [1 1 2 3 4 5 6 9] (kopya doner)
np.argsort(a)                   # siralayan INDEKSLER -> [1 3 6 0 2 7 4 5]
a.argmax(), a.argmin()          # (5, 1) en buyuk/kucuk indeksi

np.unique(a)                    # [1 2 3 4 5 6 9] tekrarsiz
np.unique(a, return_counts=True)# degerler + kac kez gectikleri

# En buyuk 3 degeri getir (argsort ile)
a[np.argsort(a)[-3:]]           # [5 6 9]

# Kosula uyan elemanlarin INDEKSLERI
np.nonzero(a > 3)               # (array([2, 4, 5, 7]),)

# Kumulatif ve istatistik
a.cumsum()                      # kumulatif toplam
np.percentile(a, [25, 50, 75])  # ceyrekler
np.isnan(np.array([1, np.nan])) # [False True] -> eksik deger bulma
\end{lstlisting}

\gsub{Neden Bu Kadar Hızlı? Bellek Düzeni}

NumPy hızının sırrı: veri bellekte \emph{tek bir bitişik blokta} durur ve
işlemler CPU'ya yakın, önbellek dostu şekilde çalışır. \code{strides}, bir
eksende bir adım atmak için kaç bayt atlanacağını söyler; \code{reshape} ve
transpoz genelde veriyi kopyalamaz, sadece bu ``adım'' bilgisini değiştirir.

\begin{lstlisting}
a = np.arange(12).reshape(3, 4)
a.strides            # (32, 8) -> satir atla=32 bayt, sutun atla=8 bayt
a.flags['C_CONTIGUOUS']  # True -> satir-oncelikli (C duzeni)

# Transpoz veriyi KOPYALAMAZ, sadece strides'i takas eder (bedava)
a.T.strides          # (8, 32)

# .ravel() bitisikse gorunum, degilse kopya doner
a.ravel()            # duzlestir (mumkunse gorunum)
np.ascontiguousarray(a.T)  # gerekince bellekte gercekten sirala
\end{lstlisting}

\ipucu{Pratik özet: (1) \textbf{Döngü yazma} --- işlemi tüm diziye uygula.
(2) \textbf{Doğru dtype seç} --- gereksiz \code{float64} bellek yer.
(3) \textbf{Görünüm/kopya farkını bil} --- beklenmedik değişiklikleri önler.
(4) Aradığın işlem muhtemelen zaten \code{np.} altında vardır --- önce ara,
sonra elle yazma.}

% #####################################################################
\gpart{2}{Pandas ile Veri İşleme}
% #####################################################################

% =====================================================================
\gsec{Pandas Temelleri: Series ve DataFrame}
% =====================================================================

Pandas, etiketli ve tablo biçimli veriyle çalışmak için tasarlanmış
kütüphanedir; adını ``panel data''dan alır. İki temel yapısı vardır:
\code{Series} (etiketli tek boyutlu dizi) ve \code{DataFrame} (etiketli
satır ve sütunlardan oluşan iki boyutlu tablo --- bir Excel sayfası gibi).
Gerçek veri biliminin büyük kısmı Pandas'ta geçer.

\gsub{Series: Etiketli Tek Boyutlu Veri}

Series, bir NumPy dizisi ile onu tanımlayan bir \emph{indeksin} (etiketler)
birleşimidir. Sözlük ve dizinin özelliklerini birleştirir.

\begin{lstlisting}
import pandas as pd

s = pd.Series([10, 20, 30, 40], index=['a', 'b', 'c', 'd'])
s['b']            # 20   -> etiketle erisim
s[1]              # 20   -> konumla erisim
s[s > 15]         # b,c,d -> boole filtreleme
s * 2             # her elemani ikiye katla (vektorlestirilmis)
s.mean()          # 25.0

# Sozlukten Series
population = pd.Series({'Ankara': 5.7, 'Izmir': 4.4, 'Istanbul': 15.8})
\end{lstlisting}

\gsub{DataFrame: Tablo Verisi}

DataFrame, her biri bir \code{Series} olan sütunların birleşimidir. Veri
biliminin ana çalışma nesnesidir.

\begin{lstlisting}
import pandas as pd

df = pd.DataFrame({
    'name':   ['Ali', 'Ayse', 'Can', 'Deniz'],
    'age':    [25, 30, 35, 28],
    'city':   ['Ankara', 'Izmir', 'Ankara', 'Istanbul'],
    'salary': [5000, 7000, 6500, 8000]
})

df.head()         # ilk 5 satir
df.tail(2)        # son 2 satir
df.shape          # (4, 4)
df.columns        # sutun adlari
df.dtypes         # her sutunun tipi
df.info()         # genel ozet (tip, bellek, eksik)
df.describe()     # sayisal sutunlarin istatistigi
\end{lstlisting}

\gsub{İlk Bakış: describe ve info}

Yeni bir veri setiyle karşılaştığında ilk yapılacak, onu ``tanımaktır''.
Bu iki metot bir veri setinin röntgenini çeker.

\begin{center}
\begin{tabularx}{\textwidth}{@{}>{\ttfamily\bfseries\color{anaMavi}}l X@{}}
\toprule
\rowcolor{acikMavi}\color{anaMavi}\textbf{Metot} & \color{anaMavi}\textbf{Ne gösterir?}\\
\midrule
df.info() & Satır sayısı, sütun tipleri, eksik olmayan değer sayısı, bellek.\\
\rowcolor{acikGri} df.describe() & Sayısal sütunlar için sayı, ortalama, std, min, çeyrekler, max.\\
df.nunique() & Her sütundaki benzersiz değer sayısı.\\
\rowcolor{acikGri} df['s'].value\_counts() & Bir sütundaki her değerin kaç kez geçtiği.\\
\bottomrule
\end{tabularx}
\end{center}

% =====================================================================
\gsec{Veri Okuma ve Yazma}
% =====================================================================

Pandas, çok çeşitli kaynaklardan veri okuyabilir. Gerçek projelerde en sık
CSV ve Excel kullanılır, ama veritabanları ve API'ler de yaygındır.

\begin{lstlisting}
import pandas as pd

# --- OKUMA ---
df = pd.read_csv('veri.csv')                  # CSV
df = pd.read_csv('veri.csv', sep=';',         # ozel ayirici
                 encoding='utf-8', na_values=['NA', '?'])
df = pd.read_excel('veri.xlsx', sheet_name='Sayfa1')
df = pd.read_json('veri.json')
df = pd.read_sql('SELECT * FROM table', connection)  # veritabani

# --- YAZMA ---
df.to_csv('cikti.csv', index=False)           # index=False: satir no yazma
df.to_excel('cikti.xlsx', index=False)
df.to_json('cikti.json', orient='records')
\end{lstlisting}

\bilgi{\code{read\_csv} onlarca parametre alır. En sık kullanılanlar:
\code{sep} (ayırıcı), \code{header} (başlık satırı), \code{index\_col}
(indeks olarak kullanılacak sütun), \code{usecols} (sadece belirli sütunlar),
\code{parse\_dates} (tarih sütunlarını otomatik dönüştür), \code{nrows}
(büyük dosyalarda ilk N satır). Büyük dosyaları \code{chunksize} ile parça
parça da okuyabilirsin.}

% =====================================================================
\gsec{Veri Seçme ve Filtreleme}
% =====================================================================

DataFrame'den veri seçmenin birden çok yolu vardır. En önemli ayrım:
\code{loc} (etiket tabanlı) ve \code{iloc} (konum/tam sayı tabanlı).

\gsub{Sütun ve Satır Seçme}

\begin{lstlisting}
df['age']                  # tek sutun -> Series
df[['name', 'age']]        # birden cok sutun -> DataFrame

# loc: ETIKET tabanli (satir etiketi, sutun adi)
df.loc[0]                  # 0 etiketli satir
df.loc[0, 'name']          # tek hucre
df.loc[0:2, ['name', 'age']]  # satir 0-2 (DAHIL), belirli sutunlar

# iloc: KONUM tabanli (tam sayi indeksleri)
df.iloc[0]                 # ilk satir
df.iloc[0:2]               # ilk 2 satir (2 HARIC, Python dilimi gibi)
df.iloc[-1]                # son satir
df.iloc[0:2, 0:2]          # ilk 2 satir, ilk 2 sutun
\end{lstlisting}

\uyari{\code{loc} ile dilimlemede son sınır \textbf{dahildir}
(\code{df.loc[0:2]} $\to$ 0,1,2), ama \code{iloc} ve normal Python
dilimlemesinde son sınır \textbf{hariçtir} (\code{df.iloc[0:2]} $\to$ 0,1).
Bu, yeni başlayanların en sık yaptığı hatalardan biridir.}

\gsub{Boole Filtreleme (En Önemli Beceri)}

Koşullara göre satır seçmek veri analizinin bel kemiğidir. Koşullar boole
maskesi üretir; bu maske DataFrame'i filtreler.

\begin{lstlisting}
# Tek kosul
df[df['age'] > 28]

# Birden cok kosul: & (ve), | (veya), ~ (degil) -- PARANTEZ sart!
df[(df['age'] > 25) & (df['city'] == 'Ankara')]
df[(df['salary'] > 6000) | (df['age'] < 26)]

# isin: bir listedeki degerler
df[df['city'].isin(['Ankara', 'Izmir'])]

# between: aralik
df[df['age'].between(25, 32)]

# String iceren
df[df['name'].str.startswith('A')]
df[df['name'].str.contains('e')]

# query: SQL benzeri okunabilir sozdizimi
df.query('age > 28 and city == "Ankara"')
\end{lstlisting}

\gsub{Yeni Sütun Oluşturma ve Dönüştürme}

\begin{lstlisting}
df['annual_salary'] = df['salary'] * 12       # vektorlestirilmis islem
df['age_group'] = df['age'].apply(
    lambda a: 'young' if a < 30 else 'adult')

# np.where ile kosullu sutun
import numpy as np
df['level'] = np.where(df['salary'] > 6500, 'high', 'normal')

# map ile deger esleme
df['city_code'] = df['city'].map({'Ankara': 6, 'Izmir': 35, 'Istanbul': 34})

# Sutun/satir silme
df = df.drop(columns=['city_code'])
df = df.drop(index=[0])
\end{lstlisting}

% =====================================================================
\gsec{Veri Temizleme}
% =====================================================================

Gerçek veri asla temiz gelmez: eksik değerler, tekrarlar, yanlış tipler ve
tutarsızlıklar içerir. Zamanının çoğunu buna harcarsın. Temiz veri, iyi
analizin ön koşuludur.

\gsub{Eksik Veri (NaN) ile Başa Çıkma}

\begin{lstlisting}
# Eksikleri tespit et
df.isnull().sum()          # her sutundaki eksik sayisi
df.isnull().sum().sum()    # toplam eksik

# Eksik satirlari/sutunlari SIL
df.dropna()                # eksik iceren satirlari at
df.dropna(axis=1)          # eksik iceren sutunlari at
df.dropna(subset=['salary']) # sadece belirli sutunda eksik olanlari at
df.dropna(thresh=3)        # en az 3 dolu deger olan satirlari tut

# Eksikleri DOLDUR (imputation)
df['salary'].fillna(0)                        # sabitle
df['salary'].fillna(df['salary'].mean())      # ortalama ile (sik)
df['salary'].fillna(df['salary'].median())    # medyan ile (aykiri degere dayanikli)
df['city'].fillna(df['city'].mode()[0])   # en sik deger (kategorik)
df.fillna(method='ffill')                 # onceki degerle (zaman serisi)
\end{lstlisting}

\notk[Silmek mi doldurmak mı?]{Eksik veri azsa (\%5'ten az) satırları silmek
güvenlidir. Ama çoksa silmek veriyi ziyan eder --- doldurmak (imputation)
gerekir. Sayısal sütunlarda \emph{medyan} (aykırı değerlere dayanıklı),
kategoriklerde \emph{mod} sık kullanılır. Neden eksik olduklarını da
düşün: rastgele mi, yoksa sistematik mi?}

\gsub{Tekrarlar ve Tip Dönüşümleri}

\begin{lstlisting}
# Tekrarlar
df.duplicated().sum()      # kac tekrar satir var
df = df.drop_duplicates()  # tekrarlari sil
df = df.drop_duplicates(subset=['name'])   # belirli sutuna gore

# Tip donusumu
df['age'] = df['age'].astype(int)
df['date'] = pd.to_datetime(df['date'])
df['salary'] = pd.to_numeric(df['salary'], errors='coerce')  # hatalari NaN yap

# String temizleme
df['name'] = df['name'].str.strip()        # bosluklari kirp
df['city'] = df['city'].str.lower()        # kucuk harf
df['name'] = df['name'].str.replace('.', '', regex=False)

# Sutun yeniden adlandirma
df = df.rename(columns={'name': 'full_name', 'age': 'age_years'})
\end{lstlisting}

% =====================================================================
\gsec{Gruplama ve Toplulaştırma (GroupBy)}
% =====================================================================

\code{groupby}, verinin belki de en güçlü aracıdır: veriyi bir veya daha çok
sütuna göre gruplara böler, her gruba bir işlem (ortalama, toplam, sayım\dots)
uygular ve sonuçları birleştirir. Buna \emph{böl-uygula-birleştir}
(split-apply-combine) denir.

\begin{lstlisting}
# Sehre gore ortalama maas
df.groupby('city')['salary'].mean()

# Birden cok istatistik
df.groupby('city')['salary'].agg(['mean', 'min', 'max', 'count'])

# Birden cok sutuna gore gruplama
df.groupby(['city', 'age_group'])['salary'].mean()

# Her sutuna farkli fonksiyon
df.groupby('city').agg({
    'salary': 'mean',
    'age':    'max',
    'name':   'count'
})

# Ozel fonksiyon
df.groupby('city')['salary'].agg(lambda x: x.max() - x.min())  # aralik
\end{lstlisting}

\gsub{Pivot Tablolar}

Pivot tablo, veriyi iki boyutlu bir özet tabloya dönüştürür --- Excel'deki
pivot tablonun aynısı. Bir sütunu satırlara, diğerini sütunlara yayar.

\begin{lstlisting}
pd.pivot_table(df,
    values='salary',      # ozetlenecek deger
    index='city',         # satirlar
    columns='age_group',  # sutunlar
    aggfunc='mean',       # toplulastirma fonksiyonu
    fill_value=0)         # bos hucreleri doldur

# crosstab: sayim icin pratik
pd.crosstab(df['city'], df['age_group'])
\end{lstlisting}

% =====================================================================
\gsec{Veri Birleştirme}
% =====================================================================

Gerçek projelerde veri genelde birden çok tablodan gelir ve birleştirilmesi
gerekir. Pandas bunun için \code{merge} (SQL join gibi), \code{concat}
(uç uca ekleme) ve \code{join} sunar.

\gsub{merge: Anahtar Üzerinden Birleştirme}

\begin{lstlisting}
employees = pd.DataFrame({
    'id': [1, 2, 3], 'name': ['Ali', 'Ayse', 'Can'], 'dept_id': [10, 20, 10]})
departments = pd.DataFrame({
    'dept_id': [10, 20, 30], 'dept': ['Software', 'Sales', 'HR']})

# Ortak sutun (dept_id) uzerinden birlestir
pd.merge(employees, departments, on='dept_id', how='inner')
\end{lstlisting}

\begin{center}
\begin{tabularx}{\textwidth}{@{}>{\ttfamily\bfseries\color{anaMavi}}l X@{}}
\toprule
\rowcolor{acikMavi}\color{anaMavi}\textbf{how} & \color{anaMavi}\textbf{Hangi satırlar tutulur?}\\
\midrule
inner & Sadece her iki tabloda da eşleşen anahtarlar (kesişim).\\
\rowcolor{acikGri} left & Sol tablonun tümü + sağdan eşleşenler.\\
right & Sağ tablonun tümü + soldan eşleşenler.\\
\rowcolor{acikGri} outer & Her iki tablonun tümü (birleşim); eşleşmeyenler NaN.\\
\bottomrule
\end{tabularx}
\end{center}

\gsub{concat: Uç Uca Ekleme}

\begin{lstlisting}
# Satir bazli birlestirme (alt alta) -- ayni sutunlar
pd.concat([df1, df2], axis=0, ignore_index=True)

# Sutun bazli birlestirme (yan yana) -- ayni satirlar
pd.concat([df1, df2], axis=1)
\end{lstlisting}

% =====================================================================
\gsec{Zaman Serileri}
% =====================================================================

Zaman serisi verisi (tarih/saat damgalı ölçümler) veri biliminde çok yaygındır:
borsa, hava, satış\dots Pandas bunun için güçlü araçlar sunar.

\begin{lstlisting}
# Tarih indeksi olustur
df['date'] = pd.to_datetime(df['date'])
df = df.set_index('date')

# Tarih bilesenlerine erisim
df.index.year
df.index.month
df.index.dayofweek         # 0=Pazartesi

# Tarihe gore dilimleme
df['2026']                 # 2026 yilinin tumu
df['2026-01':'2026-03']    # ilk uc ay

# Yeniden ornekleme (resample): gunluk -> aylik
df['sales'].resample('M').sum()       # aylik toplam
df['sales'].resample('W').mean()      # haftalik ortalama

# Kayan pencere (hareketli ortalama) -- trend icin
df['mean_7d'] = df['sales'].rolling(window=7).mean()

# Kaydirma ve degisim
df['previous'] = df['sales'].shift(1)        # bir onceki deger
df['change'] = df['sales'].pct_change()      # yuzde degisim
\end{lstlisting}

\ipucu{Zaman serilerinde \code{resample} (zaman frekansını değiştirme) ve
\code{rolling} (hareketli pencere) en çok kullanılan iki araçtır. Hareketli
ortalama, gürültülü veride \emph{trendi} görmenin en basit yoludur; günlük
dalgalanmaları yumuşatır.}

% #####################################################################
\gpart{3}{Veri Görselleştirme}
% #####################################################################

% =====================================================================
\gsec{Matplotlib: Temel Çizim}
% =====================================================================

Matplotlib, Python'ın temel çizim kütüphanesidir; diğer görselleştirme
kütüphaneleri (Seaborn, Pandas'ın çizimi) onun üzerine kuruludur. Güçlü ve
esnektir ama biraz ayrıntılıdır. İyi bir görselleştirme, bir tablodan çok daha
hızlı içgörü verir --- ``bir resim bin satıra bedeldir''.

\gsub{Anatomi: Figure ve Axes}

Matplotlib'de bir grafik iki katmandan oluşur: \code{Figure} (tüm tuval) ve
\code{Axes} (üzerindeki tek bir grafik alanı --- eksenler, veri burada çizilir).
Modern ve önerilen kullanım, bu ikisini açıkça oluşturan \emph{nesne yönelimli}
arayüzdür.

\begin{lstlisting}
import matplotlib.pyplot as plt
import numpy as np

x = np.linspace(0, 10, 100)
y = np.sin(x)

fig, ax = plt.subplots(figsize=(8, 4))    # Figure + tek Axes
ax.plot(x, y, label='sin(x)', color='navy', linewidth=2)
ax.set_title('Sinus Egrisi')
ax.set_xlabel('x')
ax.set_ylabel('sin(x)')
ax.legend()
ax.grid(True, alpha=0.3)
plt.savefig('grafik.png', dpi=150, bbox_inches='tight')  # dosyaya kaydet
plt.show()
\end{lstlisting}

\gsub{Temel Grafik Türleri}

Doğru grafik türünü seçmek, verinin doğasına bağlıdır.

\begin{center}
\begin{tabularx}{\textwidth}{@{}>{\ttfamily\bfseries\color{anaMavi}}l X@{}}
\toprule
\rowcolor{acikMavi}\color{anaMavi}\textbf{Grafik} & \color{anaMavi}\textbf{Ne zaman kullanılır?}\\
\midrule
plot & Çizgi grafik; trend, zaman serisi.\\
\rowcolor{acikGri} scatter & Dağılım; iki sayısal değişken arası ilişki.\\
bar & Çubuk; kategoriler arası karşılaştırma.\\
\rowcolor{acikGri} hist & Histogram; tek değişkenin dağılımı.\\
boxplot & Kutu grafiği; dağılım + aykırı değerler.\\
\rowcolor{acikGri} pie & Pasta; bütünün parçaları (dikkatli kullan).\\
\bottomrule
\end{tabularx}
\end{center}

\begin{lstlisting}
fig, ax = plt.subplots()

ax.scatter(df['age'], df['salary'])            # dagilim
ax.bar(df['name'], df['salary'])               # cubuk
ax.hist(df['salary'], bins=20)                 # histogram
ax.boxplot([group1, group2], labels=['A', 'B']) # kutu grafigi
ax.pie(values, labels=labels, autopct='%1.1f%%')  # pasta
\end{lstlisting}

\gsub{Birden Çok Alt Grafik (Subplots)}

Genelde birden çok grafiği yan yana göstermek isteriz. \code{plt.subplots}
bir \emph{ızgara} oluşturur.

\begin{lstlisting}
fig, axes = plt.subplots(2, 2, figsize=(10, 8))   # 2x2 izgara

axes[0, 0].plot(x, np.sin(x));  axes[0, 0].set_title('sin')
axes[0, 1].plot(x, np.cos(x));  axes[0, 1].set_title('cos')
axes[1, 0].hist(df['salary'])
axes[1, 1].scatter(df['age'], df['salary'])

fig.suptitle('Panel Grafikler', fontsize=14)
fig.tight_layout()                             # cakismalari onle
plt.show()
\end{lstlisting}

\gsub{Stil ve Özelleştirme}

\begin{lstlisting}
plt.style.use('seaborn-v0_8')                  # hazir stil temasi

ax.plot(x, y,
        color='#2E6DA4', linestyle='--', linewidth=2,
        marker='o', markersize=4, alpha=0.8, label='veri')

ax.set_xlim(0, 10)                             # eksen sinirlari
ax.set_ylim(-1.5, 1.5)
ax.set_xticks([0, 2, 4, 6, 8, 10])             # isaret konumlari
ax.axhline(0, color='gray', linewidth=0.8)     # yatay referans cizgisi
ax.annotate('tepe', xy=(1.57, 1), xytext=(3, 1.2),
            arrowprops=dict(arrowstyle='->'))  # ok ile aciklama
\end{lstlisting}

\bilgi{\textbf{pyplot vs nesne-yönelimli:} \code{plt.plot()} gibi durum tabanlı
(stateful) arayüz hızlı denemeler için pratiktir, ama karmaşık grafiklerde
kafa karıştırır. Ciddi işte her zaman \code{fig, ax = plt.subplots()} ile
başlayıp \code{ax} üzerinden çalış --- daha okunaklı ve kontrol edilebilirdir.}

% =====================================================================
\gsec{Matplotlib Derinlemesine: Grafik Türleri Kataloğu}
% =====================================================================

Yukarıdaki temeller grafik iskeletini kurar. Şimdi her grafik türünü
\emph{çalışan, tam örneklerle} inceleyelim. Her örnek kendi başına çalışır ---
kopyala, çalıştır, değiştirerek öğren.

\gsub{Çizgi Grafik: Birden Çok Seri ve Stil}

\begin{lstlisting}
import matplotlib.pyplot as plt
import numpy as np

x = np.linspace(0, 10, 200)
fig, ax = plt.subplots(figsize=(8, 4))

ax.plot(x, np.sin(x), color='#2E6DA4', lw=2, label='sin')
ax.plot(x, np.cos(x), color='#D9701E', lw=2, ls='--', label='cos')
ax.plot(x, np.sin(x)*np.exp(-x/5), color='#2E7D32', lw=2,
        ls=':', label='sonumlu')

ax.fill_between(x, np.sin(x), 0, alpha=0.1, color='#2E6DA4')  # alan doldur
ax.set(title='Cizgi Grafikleri', xlabel='x', ylabel='y')
ax.legend(loc='upper right', frameon=True)
ax.grid(True, alpha=0.3)
plt.show()
\end{lstlisting}

\gsub{Dağılım Grafiği: Renk ve Boyutla Kodlama}

Bir dağılım grafiğinde 4 boyutu aynı anda gösterebilirsin: x, y, renk (\code{c})
ve nokta boyutu (\code{s}).

\begin{lstlisting}
rng = np.random.default_rng(0)
x = rng.random(80); y = rng.random(80)
renk = rng.random(80)              # 3. boyut -> renk
boyut = rng.random(80) * 400       # 4. boyut -> alan

fig, ax = plt.subplots(figsize=(7, 5))
sc = ax.scatter(x, y, c=renk, s=boyut, cmap='viridis',
                alpha=0.7, edgecolors='white', linewidth=0.5)
fig.colorbar(sc, ax=ax, label='renk degeri')  # renk olcegi
ax.set(title='Dagilim: 4 boyut', xlabel='x', ylabel='y')
plt.show()
\end{lstlisting}

\gsub{Çubuk Grafik: Gruplu ve Yığılmış}

\begin{lstlisting}
kategoriler = ['Ocak', 'Subat', 'Mart', 'Nisan']
urun_a = [20, 34, 30, 35]
urun_b = [25, 32, 34, 20]
x = np.arange(len(kategoriler))
w = 0.35

fig, (ax1, ax2) = plt.subplots(1, 2, figsize=(11, 4))

# Gruplu (yan yana)
ax1.bar(x - w/2, urun_a, w, label='A', color='#2E6DA4')
ax1.bar(x + w/2, urun_b, w, label='B', color='#D9701E')
ax1.set_xticks(x); ax1.set_xticklabels(kategoriler)
ax1.set_title('Gruplu Cubuk'); ax1.legend()

# Yigilmis (ust uste)
ax2.bar(kategoriler, urun_a, label='A', color='#2E6DA4')
ax2.bar(kategoriler, urun_b, bottom=urun_a, label='B', color='#D9701E')
ax2.set_title('Yigilmis Cubuk'); ax2.legend()

fig.tight_layout(); plt.show()
\end{lstlisting}

\gsub{Histogram ve Yoğunluk: Dağılımı Görmek}

\begin{lstlisting}
veri1 = np.random.default_rng(1).normal(60, 10, 1000)
veri2 = np.random.default_rng(2).normal(75, 8, 1000)

fig, ax = plt.subplots(figsize=(8, 4))
ax.hist(veri1, bins=30, alpha=0.6, color='#2E6DA4', label='Grup 1',
        density=True)                          # density: alan=1
ax.hist(veri2, bins=30, alpha=0.6, color='#D9701E', label='Grup 2',
        density=True)
ax.axvline(veri1.mean(), color='#2E6DA4', ls='--', lw=2)  # ortalama cizgisi
ax.axvline(veri2.mean(), color='#D9701E', ls='--', lw=2)
ax.set(title='Iki Grubun Dagilimi', xlabel='puan', ylabel='yogunluk')
ax.legend(); plt.show()
\end{lstlisting}

\gsub{Kutu ve Keman: Grupları Karşılaştırmak}

\begin{lstlisting}
gruplar = [np.random.default_rng(i).normal(50+i*10, 8, 100)
           for i in range(4)]

fig, (ax1, ax2) = plt.subplots(1, 2, figsize=(11, 4))
ax1.boxplot(gruplar, tick_labels=['A', 'B', 'C', 'D'])
ax1.set_title('Kutu Grafigi (medyan, ceyrek, aykiri)')

parts = ax2.violinplot(gruplar, showmeans=True)  # keman: dagilim sekli
ax2.set_xticks([1, 2, 3, 4]); ax2.set_xticklabels(['A','B','C','D'])
ax2.set_title('Keman Grafigi')
fig.tight_layout(); plt.show()
\end{lstlisting}

% =====================================================================
\gsec{Matplotlib Derinlemesine: Eksen, Yazı ve Kaydetme}
% =====================================================================

\gsub{Eksenleri Tam Kontrol: Sınırlar, İşaretler, Ölçek}

\begin{lstlisting}
import matplotlib.ticker as mticker

fig, ax = plt.subplots(figsize=(8, 4))
x = np.linspace(1, 100, 100)
ax.plot(x, x**2)

ax.set_xlim(0, 100); ax.set_ylim(1, 1e4)
ax.set_yscale('log')                       # logaritmik eksen
ax.set_xticks([0, 25, 50, 75, 100])
ax.set_xticklabels(['0', 'ceyrek', 'yari', '3/4', 'tam'], rotation=45)

# Ekseni yuzde/para birimi olarak bicimle
ax.yaxis.set_major_formatter(mticker.FuncFormatter(
    lambda v, _: f'{v:,.0f} TL'))
ax.tick_params(axis='both', labelsize=9)
plt.show()
\end{lstlisting}

\gsub{Metin, Ok ve Vurgulama (Annotation)}

Bir grafiği anlatan şey çoğu zaman üzerindeki açıklamalardır: bir zirveyi
işaretlemek, bir eşiği göstermek.

\begin{lstlisting}
x = np.linspace(0, 10, 100)
y = np.sin(x)
fig, ax = plt.subplots(figsize=(8, 4))
ax.plot(x, y, color='#2E6DA4', lw=2)

# Bir noktaya ok ile aciklama
tepe_x = np.pi/2
ax.annotate('en yuksek nokta',
            xy=(tepe_x, 1), xytext=(4, 1.3),
            arrowprops=dict(arrowstyle='->', color='#C0392B'),
            fontsize=11, color='#C0392B')

# Serbest metin ve referans cizgileri
ax.text(7, -0.8, r'$y=\sin(x)$', fontsize=12)  # LaTeX matematik
ax.axhline(0, color='gray', lw=0.8)
ax.axvspan(3, 6, alpha=0.1, color='#D9701E')   # dikey bant vurgusu
ax.set(title='Aciklamali Grafik', ylim=(-1.5, 1.6))
plt.show()
\end{lstlisting}

\gsub{İkiz Eksen (Twin Axes): Farklı Birimler}

Sıcaklık ve satış gibi farklı ölçekli iki seriyi tek grafikte göstermek için.

\begin{lstlisting}
aylar = np.arange(1, 13)
sicaklik = [5, 6, 10, 15, 20, 25, 28, 27, 22, 16, 10, 6]
satis = [200, 180, 220, 300, 450, 600, 700, 680, 500, 350, 250, 210]

fig, ax1 = plt.subplots(figsize=(8, 4))
ax1.plot(aylar, sicaklik, 'o-', color='#C0392B')
ax1.set_xlabel('ay'); ax1.set_ylabel('sicaklik (C)', color='#C0392B')
ax1.tick_params(axis='y', labelcolor='#C0392B')

ax2 = ax1.twinx()                          # ayni x, farkli y ekseni
ax2.bar(aylar, satis, alpha=0.3, color='#2E6DA4')
ax2.set_ylabel('satis (adet)', color='#2E6DA4')
ax2.tick_params(axis='y', labelcolor='#2E6DA4')
fig.suptitle('Sicaklik vs Satis')
plt.show()
\end{lstlisting}

\gsub{Karmaşık Yerleşim ve Yüksek Kaliteli Kaydetme}

\begin{lstlisting}
# Farkli boyutlarda paneller icin mosaic (en pratik yontem)
fig, axd = plt.subplot_mosaic(
    [['ust', 'ust'],
     ['sol', 'sag']], figsize=(9, 6))
axd['ust'].plot(np.random.default_rng(0).random(50).cumsum())
axd['sol'].hist(np.random.default_rng(1).normal(size=200))
axd['sag'].scatter(np.random.default_rng(2).random(50),
                   np.random.default_rng(3).random(50))
for ad, ax in axd.items():
    ax.set_title(ad)
fig.tight_layout()

# Yayina uygun kaydetme
fig.savefig('rapor.png', dpi=300, bbox_inches='tight')  # yuksek cozunurluk
fig.savefig('rapor.pdf', bbox_inches='tight')           # vektorel (olcek bozulmaz)
fig.savefig('rapor.svg')                                 # web/duzenlenebilir
\end{lstlisting}

\bilgi{\textbf{Kaydetme ipuçları:} Raporlar için \code{dpi=300} kullan; sunum ve
web için \code{bbox\_inches='tight'} kenar boşluklarını kırpar. Ölçeklendiğinde
bulanıklaşmaması gereken grafikleri (dergi, poster) \code{.pdf} veya \code{.svg}
olarak vektörel kaydet --- \code{.png} pikseldir, büyütünce bozulur.}

\ipucu{Matplotlib'i ezberlemeye çalışma; mantığını kavra: her şey bir
\code{Figure} (tuval) ve bir veya çok \code{Axes} (grafik alanı) üzerinedir. Bir
grafik türü lazım olduğunda ``matplotlib \emph{[grafik adı}] example'' diye ara
--- galeri (\href{https://matplotlib.org/stable/gallery/}%
{matplotlib.org/stable/gallery}) neredeyse her ihtiyaca hazır, kopyalanabilir
örnek sunar.}

% =====================================================================
\gsec{Pandas ile Hızlı Çizim}
% =====================================================================

Pandas, keşif sırasında hızlı grafikler için Matplotlib'i doğrudan sarar.
Bir DataFrame veya Series üzerinde \code{.plot()} çağırman yeterlidir.

\begin{lstlisting}
df['salary'].plot(kind='hist', bins=20, title='Maas Dagilimi')
df['salary'].plot(kind='box')
df.plot(x='age', y='salary', kind='scatter')
df.groupby('city')['salary'].mean().plot(kind='bar')
df['sales'].plot(kind='line')                  # zaman serisi

# kind secenekleri: 'line','bar','barh','hist','box','kde','area','scatter','pie'
\end{lstlisting}

% =====================================================================
\gsec{Seaborn: İstatistiksel Görselleştirme}
% =====================================================================

Seaborn, Matplotlib üzerine kurulu, istatistiksel grafikler için tasarlanmış
bir kütüphanedir. Daha az kodla daha güzel ve bilgilendirici grafikler üretir;
DataFrame'lerle doğrudan çalışır ve kategorik değişkenleri anlar.

\gsub{Dağılım Grafikleri}

\begin{lstlisting}
import seaborn as sns
import matplotlib.pyplot as plt

sns.set_theme(style='whitegrid')               # sik varsayilan tema

# Histogram + yogunluk egrisi
sns.histplot(data=df, x='salary', kde=True, bins=20)

# Yogunluk (KDE) grafigi
sns.kdeplot(data=df, x='salary', hue='city', fill=True)

# Kutu ve keman grafikleri (grup karsilastirma)
sns.boxplot(data=df, x='city', y='salary')
sns.violinplot(data=df, x='city', y='salary')  # kutu + yogunluk

plt.show()
\end{lstlisting}

\gsub{İlişki Grafikleri}

İki değişken arasındaki ilişkiyi görselleştirmek, korelasyon ve modelleme için
kritiktir.

\begin{lstlisting}
# Dagilim + regresyon dogrusu
sns.scatterplot(data=df, x='age', y='salary', hue='city', size='salary')
sns.regplot(data=df, x='age', y='salary')      # dogrusal uyum + guven araligi

# Pairplot: TUM sayisal sutun ciftleri (kesif icin harika)
sns.pairplot(df, hue='city')

# Isi haritasi (heatmap) -- korelasyon matrisi icin ideal
correlation = df.corr(numeric_only=True)
sns.heatmap(correlation, annot=True, cmap='coolwarm', center=0)
\end{lstlisting}

\gsub{Kategorik Grafikler}

\begin{lstlisting}
sns.barplot(data=df, x='city', y='salary')     # ortalama + guven araligi
sns.countplot(data=df, x='city')               # kategori sayimi
sns.catplot(data=df, x='city', y='salary',     # panel katmanli
            kind='box', col='age_group')
\end{lstlisting}

\ipucu{Keşifsel analiz (EDA) için Seaborn'un \code{pairplot} ve \code{heatmap}
fonksiyonları paha biçilmezdir: \code{pairplot} tüm değişken çiftlerinin
ilişkisini tek bakışta gösterir; korelasyon \code{heatmap}'i hangi
değişkenlerin birbiriyle ilişkili olduğunu (ve modelde çoklu bağlantı riskini)
anında ortaya koyar. Yeni bir veri setinde ilk çalıştıracağın şeyler bunlar
olmalı.}

% #####################################################################
\gpart{4}{İstatistik ve Olasılık}
% #####################################################################

% =====================================================================
\gsec{Betimsel İstatistik}
% =====================================================================

İstatistik, veri biliminin teorik temelidir. Betimsel istatistik, bir veri
setini birkaç sayıyla \emph{özetler}: merkezi eğilim (veri nerede toplanıyor?)
ve yayılım (ne kadar dağınık?). Bu özetler, veriyi anlamanın ilk adımıdır.

\gsub{Merkezi Eğilim: Ortalama, Medyan, Mod}

\begin{center}
\begin{tabularx}{\textwidth}{@{}>{\bfseries\color{anaMavi}}l X X@{}}
\toprule
\rowcolor{acikMavi}\color{anaMavi}\textbf{Ölçü} & \color{anaMavi}\textbf{Tanım} & \color{anaMavi}\textbf{Özellik}\\
\midrule
Ortalama (mean) & Tüm değerlerin toplamı / sayısı & Aykırı değerlerden \emph{etkilenir}\\
\rowcolor{acikGri} Medyan (median) & Sıralı verinin ortadaki değeri & Aykırı değere \emph{dayanıklı}\\
Mod (mode) & En sık geçen değer & Kategorik veride kullanışlı\\
\bottomrule
\end{tabularx}
\end{center}

\begin{lstlisting}
import numpy as np
data = np.array([10, 12, 12, 13, 15, 18, 100])   # 100 bir aykiri deger

np.mean(data)      # 25.7  -> aykiri deger yukari cekti!
np.median(data)    # 13.0  -> daha temsili
from scipy import stats
stats.mode(data, keepdims=False).mode   # 12

# Pandas ile
df['salary'].mean(); df['salary'].median(); df['salary'].mode()
\end{lstlisting}

\notk[Ortalama mı medyan mı?]{Veri \emph{simetrikse} ortalama ve medyan
birbirine yakındır. Ama veri \emph{çarpıksa} (skewed --- örneğin gelir,
ev fiyatları) veya aykırı değerler varsa, medyan daha temsilidir. ``Ortalama
maaş'' bir CEO yüzünden yanıltıcı olabilir; ``medyan maaş'' gerçeği daha iyi
yansıtır.}

\gsub{Yayılım: Varyans, Standart Sapma, Çeyrekler}

Yayılım, verinin ne kadar ``dağınık'' olduğunu ölçer. İki veri seti aynı
ortalamaya sahip olup çok farklı yayılıma sahip olabilir.

\begin{lstlisting}
np.var(data)       # varyans: ortalamadan sapmalarin karesinin ortalamasi
np.std(data)       # standart sapma: varyansin karekoku (ayni birimde)
np.ptp(data)       # aralik (range): max - min

# Ceyrekler ve IQR (ceyrekler arasi aralik)
q1 = np.percentile(data, 25)      # 1. ceyrek
q3 = np.percentile(data, 75)      # 3. ceyrek
iqr = q3 - q1                     # orta %50'nin yayilimi

# Aykiri deger tespiti (IQR yontemi)
lower_bound = q1 - 1.5 * iqr
upper_bound = q3 + 1.5 * iqr
outliers = data[(data < lower_bound) | (data > upper_bound)]
\end{lstlisting}

\uyari{NumPy varsayılan olarak \emph{anakütle} (population) varyansını hesaplar
(N'e böler). \emph{Örneklem} (sample) varyansı için (N-1'e bölme, yansız
tahmin) \code{np.var(data, ddof=1)} kullan. Pandas ise varsayılan olarak
örneklem varyansını (ddof=1) kullanır --- bu fark küçük örneklemlerde önemlidir.}

% =====================================================================
\gsec{Olasılık Dağılımları}
% =====================================================================

Olasılık dağılımı, bir rastgele değişkenin hangi değerleri hangi olasılıkla
aldığını tarif eder. Veri bilimi modellerinin çoğu, verinin belirli bir
dağılımdan geldiğini varsayar. En önemlisi \emph{normal dağılımdır}.

\gsub{Normal (Gauss) Dağılımı}

Normal dağılım (çan eğrisi), doğada ve veride en sık karşılaşılan dağılımdır:
boy, ölçüm hataları, birçok doğal olgu buna yakındır. İki parametresi vardır:
ortalama $\mu$ (merkez) ve standart sapma $\sigma$ (genişlik).

\begin{lstlisting}
import numpy as np
from scipy import stats
import matplotlib.pyplot as plt

# Normal dagilimdan orneklem uret
samples = np.random.normal(loc=170, scale=10, size=1000)  # mu=170, sigma=10

# Olasilik yogunluk fonksiyonu (PDF)
x = np.linspace(140, 200, 100)
pdf = stats.norm.pdf(x, loc=170, scale=10)
plt.plot(x, pdf)
plt.hist(samples, bins=30, density=True, alpha=0.5)

# Kumulatif olasilik: boyu 180'den kucuk olma olasiligi
stats.norm.cdf(180, loc=170, scale=10)   # ~0.84
\end{lstlisting}

\bilgi{\textbf{68-95-99.7 kuralı:} Normal dağılımda verinin \%68'i
ortalamadan $\pm 1\sigma$, \%95'i $\pm 2\sigma$, \%99.7'si $\pm 3\sigma$
içindedir. Bu yüzden ``$3\sigma$ dışı'' bir gözlem çok nadirdir (\%0.3) ve
genelde aykırı değer sayılır.}

\gsub{Diğer Önemli Dağılımlar}

\begin{center}
\begin{tabularx}{\textwidth}{@{}>{\bfseries\color{anaMavi}}l X@{}}
\toprule
\rowcolor{acikMavi}\color{anaMavi}\textbf{Dağılım} & \color{anaMavi}\textbf{Ne modeller?}\\
\midrule
Normal & Sürekli, simetrik olgular (boy, hata).\\
\rowcolor{acikGri} Binom & N denemede başarı sayısı (yazı-tura).\\
Poisson & Sabit aralıkta olay sayısı (saatte gelen çağrı).\\
\rowcolor{acikGri} Üstel (exponential) & Olaylar arası bekleme süresi.\\
Düzgün (uniform) & Tüm sonuçların eşit olasılıklı olduğu durum.\\
\bottomrule
\end{tabularx}
\end{center}

\gsub{Merkezi Limit Teoremi}

Merkezi Limit Teoremi (MLT), istatistiğin en güçlü sonucudur: \emph{hangi
dağılımdan gelirse gelsin}, yeterince büyük örneklemlerin \textbf{ortalamaları}
normal dağılıma yaklaşır. Bu, birçok istatistiksel testin neden çalıştığını
açıklar.

\begin{lstlisting}
# Kaynak dagilim: tamamen carpik (ustel) -- normal DEGIL
means = []
for _ in range(1000):
    sample = np.random.exponential(scale=2, size=50)  # 50'lik ornek
    means.append(sample.mean())                       # ortalamasini al

# means'in histogrami -> CAN EGRISINE benzer! (MLT)
plt.hist(means, bins=30)
\end{lstlisting}

% =====================================================================
\gsec{Korelasyon ve Hipotez Testleri}
% =====================================================================

\gsub{Korelasyon: İlişki Ölçümü}

Korelasyon, iki değişkenin birlikte nasıl değiştiğini $-1$ ile $+1$ arasında
ölçer. $+1$ mükemmel pozitif, $-1$ mükemmel negatif, $0$ ilişki yok demektir.

\begin{lstlisting}
# Pearson korelasyonu (dogrusal iliski)
df['age'].corr(df['salary'])

# Tum sayisal sutunlarin korelasyon matrisi
df.corr(numeric_only=True)

from scipy import stats
r, p = stats.pearsonr(df['age'], df['salary'])   # katsayi + p-degeri
rho, p = stats.spearmanr(df['age'], df['salary'])# siralamaya dayali (dogrusal olmayan)
\end{lstlisting}

\uyari{\textbf{Korelasyon nedensellik değildir!} İki değişken güçlü ilişkili
olabilir ama biri diğerine \emph{neden} olmuyor olabilir --- ikisini birden
etkileyen üçüncü bir faktör (confounder) olabilir. Klasik örnek: dondurma
satışı ve boğulma vakaları ilişkilidir, ama sebep ikisini de artıran
sıcak havadır.}

\gsub{Hipotez Testleri}

Hipotez testi, bir iddianın (örneğin ``iki grubun ortalaması farklı'')
veriyle desteklenip desteklenmediğini istatistiksel olarak sorgular. Mantık:
bir \emph{sıfır hipotezi} ($H_0$: fark yok) kurulur; veri bu hipotez altında
çok olası değilse ($p$-değeri küçükse) $H_0$ reddedilir.

\begin{lstlisting}
from scipy import stats

# t-testi: iki grubun ortalamasi farkli mi?
group_a = df[df['city'] == 'Ankara']['salary']
group_b = df[df['city'] == 'Izmir']['salary']
t_stat, p = stats.ttest_ind(group_a, group_b)

if p < 0.05:
    print("Fark istatistiksel olarak anlamli (H0 reddedildi)")
else:
    print("Anlamli fark yok")

# Ki-kare testi: iki kategorik degisken bagimsiz mi?
table = pd.crosstab(df['city'], df['age_group'])
chi2, p, dof, expected = stats.chi2_contingency(table)

# ANOVA: ikiden fazla grubun ortalamasi farkli mi?
f_stat, p = stats.f_oneway(group1, group2, group3)
\end{lstlisting}

\notk[p-değeri nedir?]{$p$-değeri, ``$H_0$ doğruysa, gözlemlediğimiz kadar
(veya daha) uç bir sonucu görme olasılığıdır''. Küçük $p$ ($< 0.05$ yaygın
eşik), sonucun şans eseri olma ihtimalinin düşük olduğunu, yani $H_0$'a
karşı kanıt olduğunu gösterir. Ama $p < 0.05$ ``kesin doğru'' demek değildir;
sadece bir kanıt eşiğidir ve dikkatli yorumlanmalıdır.}

% #####################################################################
\gpart{5}{Makine Öğrenmesi (scikit-learn)}
% #####################################################################

% =====================================================================
\gsec{Makine Öğrenmesine Giriş}
% =====================================================================

Makine öğrenmesi (ML), açıkça programlamak yerine \emph{veriden örüntü
öğrenen} algoritmalar bütünüdür. Bir model, geçmiş veriye ``bakarak'' bir
görevde (tahmin, sınıflandırma) ustalaşır. scikit-learn, klasik ML için
Python'ın standart kütüphanesidir ve tutarlı bir arayüz sunar.

\gsub{ML Türleri}

\begin{center}
\begin{tabularx}{\textwidth}{@{}>{\bfseries\color{anaMavi}}l X X@{}}
\toprule
\rowcolor{acikMavi}\color{anaMavi}\textbf{Tür} & \color{anaMavi}\textbf{Tanım} & \color{anaMavi}\textbf{Örnek}\\
\midrule
Denetimli (regresyon) & Etiketli veriden sürekli değer tahmini & Ev fiyatı\\
\rowcolor{acikGri} Denetimli (sınıflandırma) & Etiketli veriden kategori tahmini & Spam mı değil mi\\
Denetimsiz (kümeleme) & Etiketsiz veride grup bulma & Müşteri segmentleri\\
\rowcolor{acikGri} Boyut indirgeme & Değişken sayısını azaltma & PCA ile sıkıştırma\\
\bottomrule
\end{tabularx}
\end{center}

\gsub{scikit-learn'in Tutarlı Arayüzü}

scikit-learn'ün gücü tutarlılığındadır: her model aynı üç metodu paylaşır.
Bir kez öğrenince tüm modelleri kullanabilirsin.

\begin{center}
\begin{tabularx}{\textwidth}{@{}>{\ttfamily\bfseries\color{anaMavi}}l X@{}}
\toprule
\rowcolor{acikMavi}\color{anaMavi}\textbf{Metot} & \color{anaMavi}\textbf{Ne yapar?}\\
\midrule
.fit(X, y) & Modeli eğitim verisiyle eğitir (öğrenir).\\
\rowcolor{acikGri} .predict(X) & Yeni veri için tahmin üretir.\\
.transform(X) & Veriyi dönüştürür (ölçekleme, kodlama vb.).\\
\rowcolor{acikGri} .score(X, y) & Modelin başarısını ölçer.\\
\bottomrule
\end{tabularx}
\end{center}

\gsub{Eğitim/Test Ayrımı: Altın Kural}

Bir modeli, eğitildiği veriyle değerlendirmek yanıltıcıdır --- model o veriyi
``ezberlemiş'' olabilir. Bu yüzden veriyi \textbf{eğitim} ve \textbf{test}
olarak ayırırız: model eğitimde öğrenir, hiç görmediği test verisinde ölçülür.

\begin{lstlisting}
from sklearn.model_selection import train_test_split

X = df[['age', 'experience', 'education']]   # ozellikler (girdi)
y = df['salary']                             # hedef (cikti)

X_train, X_test, y_train, y_test = train_test_split(
    X, y,
    test_size=0.2,        # %20 test, %80 egitim
    random_state=42)      # tekrarlanabilirlik icin
\end{lstlisting}

\uyari{\textbf{Veri sızıntısından (data leakage) kaçın:} Ölçekleme,
eksik doldurma gibi ön işlemleri \emph{sadece eğitim verisine} \code{fit} et,
sonra hem eğitim hem teste \code{transform} uygula. Tüm veriye birden fit
edersen, test bilgisi eğitime ``sızar'' ve başarını gerçekçi olmayan biçimde
yüksek gösterir. Pipeline'lar (ileride) bunu otomatik önler.}

% =====================================================================
\gsec{Regresyon}
% =====================================================================

Regresyon, sürekli bir sayısal değeri tahmin eder (fiyat, sıcaklık, maaş).
En temel modeli \emph{doğrusal regresyondur}: çıktının, girdilerin ağırlıklı
toplamı olduğunu varsayar: $y = w_0 + w_1 x_1 + \dots + w_n x_n$.

\gsub{Doğrusal Regresyon}

\begin{lstlisting}
from sklearn.linear_model import LinearRegression
from sklearn.metrics import mean_squared_error, r2_score
import numpy as np

model = LinearRegression()
model.fit(X_train, y_train)           # katsayilari ogren

y_pred = model.predict(X_test)        # test verisinde tahmin et

# Ogrenilen parametreler
model.coef_          # her ozelligin agirligi (w1..wn)
model.intercept_     # sabit terim (w0)

# Degerlendirme
rmse = np.sqrt(mean_squared_error(y_test, y_pred))  # ort. hata (birimde)
r2 = r2_score(y_test, y_pred)         # aciklanan varyans orani (0-1)
print(f"RMSE: {rmse:.2f}, R^2: {r2:.3f}")
\end{lstlisting}

\gsub{Regresyon Metrikleri}

\begin{center}
\begin{tabularx}{\textwidth}{@{}>{\bfseries\color{anaMavi}}l X@{}}
\toprule
\rowcolor{acikMavi}\color{anaMavi}\textbf{Metrik} & \color{anaMavi}\textbf{Anlamı}\\
\midrule
MAE & Ortalama mutlak hata; yorumlaması kolay, aykırıya dayanıklı.\\
\rowcolor{acikGri} MSE / RMSE & Kare hata; büyük hataları daha çok cezalandırır.\\
$R^2$ & Modelin açıkladığı varyans oranı; 1'e yakın = iyi, 0 = ortalama kadar.\\
\bottomrule
\end{tabularx}
\end{center}

\gsub{Aşırı Öğrenme ve Düzenlileştirme (Regularization)}

Model eğitim verisini \emph{ezberlerse} (aşırı öğrenme --- overfitting), yeni
veride kötü çalışır. Ridge ve Lasso, katsayıları küçük tutarak
(düzenlileştirme) bunu önler. Lasso ayrıca bazı katsayıları sıfırlayıp
otomatik özellik seçimi yapar.

\begin{lstlisting}
from sklearn.linear_model import Ridge, Lasso

ridge = Ridge(alpha=1.0)    # L2 duzenlilestirme (katsayilari kucultur)
lasso = Lasso(alpha=0.1)    # L1 (bazi katsayilari SIFIRLAR)
ridge.fit(X_train, y_train)

# Dogrusal olmayan iliskiler icin agac tabanli:
from sklearn.ensemble import RandomForestRegressor
rf = RandomForestRegressor(n_estimators=100, random_state=42)
rf.fit(X_train, y_train)
rf.feature_importances_     # hangi ozellik ne kadar onemli
\end{lstlisting}

\bilgi{\textbf{Aşırı vs eksik öğrenme:} Eğitimde iyi, testte kötü = aşırı
öğrenme (model çok karmaşık). Hem eğitimde hem testte kötü = eksik öğrenme
(model çok basit). İkisinin dengesi ``bias-variance ödünleşimi''dir. Eğitim
ve test skorlarını \emph{birlikte} izlemek, hangisinde olduğunu gösterir.}

% =====================================================================
\gsec{Sınıflandırma}
% =====================================================================

Sınıflandırma, bir örneğin hangi \emph{kategoriye} ait olduğunu tahmin eder:
spam/değil, hasta/sağlıklı, hangi rakam. Çıktı süreklinin aksine ayrık
sınıflardır.

\gsub{Lojistik Regresyon ve Diğer Modeller}

Adına rağmen lojistik regresyon bir \emph{sınıflandırma} algoritmasıdır:
bir örneğin bir sınıfa ait olma \emph{olasılığını} tahmin eder.

\begin{lstlisting}
from sklearn.linear_model import LogisticRegression
from sklearn.tree import DecisionTreeClassifier
from sklearn.ensemble import RandomForestClassifier
from sklearn.svm import SVC
from sklearn.neighbors import KNeighborsClassifier

model = LogisticRegression(max_iter=1000)
model.fit(X_train, y_train)

y_pred = model.predict(X_test)              # sinif tahmini
probabilities = model.predict_proba(X_test) # her sinifin olasiligi

# Diger siniflandiricilar -- ayni arayuz!
tree = DecisionTreeClassifier(max_depth=5)
forest = RandomForestClassifier(n_estimators=100)
knn = KNeighborsClassifier(n_neighbors=5)
\end{lstlisting}

\gsub{Sınıflandırma Metrikleri ve Karmaşıklık Matrisi}

Doğruluk (accuracy) tek başına yanıltıcı olabilir --- özellikle
\emph{dengesiz} veride (örneğin \%99 sağlıklı). Bu yüzden precision, recall
ve F1 kullanılır.

\begin{lstlisting}
from sklearn.metrics import (accuracy_score, precision_score, recall_score,
                             f1_score, confusion_matrix, classification_report)

accuracy_score(y_test, y_pred)     # dogru tahmin orani
precision_score(y_test, y_pred)    # pozitif dedigimden kaci gercekten pozitif
recall_score(y_test, y_pred)       # gercek pozitiflerin kacini yakaladim
f1_score(y_test, y_pred)           # precision ve recall'un harmonik ortalamasi

print(confusion_matrix(y_test, y_pred))    # TP, FP, TN, FN tablosu
print(classification_report(y_test, y_pred))  # hepsi bir arada
\end{lstlisting}

\begin{center}
\begin{tabularx}{\textwidth}{@{}>{\bfseries\color{anaMavi}}l X@{}}
\toprule
\rowcolor{acikMavi}\color{anaMavi}\textbf{Metrik} & \color{anaMavi}\textbf{Ne zaman önemli?}\\
\midrule
Precision (kesinlik) & Yanlış alarm pahalıysa (spam olmayanı spam demek).\\
\rowcolor{acikGri} Recall (duyarlılık) & Kaçırmak pahalıysa (hasta olanı sağlıklı demek).\\
F1 & Precision ve recall dengesi gerektiğinde.\\
\bottomrule
\end{tabularx}
\end{center}

\notk[Precision mı recall mı?]{Bu, hatanın maliyetine bağlıdır. Kanser
teşhisinde \emph{recall} kritiktir: bir hastayı kaçırmak, yanlış alarmdan
çok daha kötüdür. Spam filtresinde \emph{precision} önemlidir: önemli bir
maili spam'e atmak kötüdür. Hangisini optimize edeceğine problem karar verir.}

% =====================================================================
\gsec{Denetimsiz Öğrenme: Kümeleme ve PCA}
% =====================================================================

Denetimsiz öğrenmede \emph{etiket yoktur}; amaç veride gizli yapı bulmaktır.
En yaygın iki görev: kümeleme (benzer örnekleri gruplama) ve boyut indirgeme.

\gsub{K-Means Kümeleme}

K-Means, veriyi $k$ kümeye ayırır: her örneği en yakın küme merkezine atar,
merkezleri günceller, tekrarlar. Müşteri segmentasyonu, görüntü sıkıştırma
kullanır.

\begin{lstlisting}
from sklearn.cluster import KMeans

model = KMeans(n_clusters=3, random_state=42, n_init=10)
clusters = model.fit_predict(X)    # her ornegin kume etiketi (0,1,2)
model.cluster_centers_             # kume merkezleri

# Kac kume? Dirsek (elbow) yontemi
inertia = []
for k in range(1, 10):
    km = KMeans(n_clusters=k, n_init=10, random_state=42).fit(X)
    inertia.append(km.inertia_)    # kume ici toplam uzaklik
# inertia grafiginde 'dirsek' noktasi ideal k'yi verir
\end{lstlisting}

\gsub{PCA: Boyut İndirgeme}

Temel Bileşen Analizi (PCA), çok sayıda değişkeni, bilginin çoğunu koruyarak
az sayıda ``bileşene'' indirger. Görselleştirme (yüksek boyutu 2B'ye indirme)
ve gürültü azaltma için kullanılır.

\begin{lstlisting}
from sklearn.decomposition import PCA
from sklearn.preprocessing import StandardScaler

# PCA olcege duyarli -> ONCE standartlastir
X_scaled = StandardScaler().fit_transform(X)

pca = PCA(n_components=2)           # 2 bilesene indir
X_pca = pca.fit_transform(X_scaled)
pca.explained_variance_ratio_      # her bilesenin acikladigi varyans

# 2B'ye indirip gorsellestir
plt.scatter(X_pca[:, 0], X_pca[:, 1], c=clusters)
\end{lstlisting}

% =====================================================================
\gsec{Model Değerlendirme ve Pipeline'lar}
% =====================================================================

\gsub{Çapraz Doğrulama (Cross-Validation)}

Tek bir eğitim/test ayrımı şansa bağlı olabilir. Çapraz doğrulama, veriyi $k$
parçaya böler; her seferinde birini test, gerisini eğitim yapıp $k$ kez ölçer
ve ortalar. Daha güvenilir bir başarı tahmini verir.

\begin{lstlisting}
from sklearn.model_selection import cross_val_score

scores = cross_val_score(model, X, y, cv=5, scoring='accuracy')  # 5-katli
print(f"Ortalama: {scores.mean():.3f} (+/- {scores.std():.3f})")
\end{lstlisting}

\gsub{Hiperparametre Ayarı: GridSearch}

Modellerin, veriden öğrenilmeyen ama elle ayarlanan \emph{hiperparametreleri}
vardır (ağaç derinliği, komşu sayısı\dots). GridSearch, bir parametre ızgarasını
sistematik deneyip en iyisini çapraz doğrulamayla bulur.

\begin{lstlisting}
from sklearn.model_selection import GridSearchCV

param_grid = {
    'n_estimators': [50, 100, 200],
    'max_depth':    [5, 10, None]
}
search = GridSearchCV(RandomForestClassifier(random_state=42),
                      param_grid, cv=5, scoring='f1')
search.fit(X_train, y_train)

search.best_params_    # en iyi parametre kombinasyonu
search.best_score_     # en iyi capraz dogrulama skoru
best_model = search.best_estimator_
\end{lstlisting}

\gsub{Pipeline: Ön İşleme + Model Tek Zincirde}

Pipeline, ön işleme adımlarını ve modeli tek bir nesnede birleştirir. Bu, hem
kodu temizler hem de veri sızıntısını otomatik önler (ön işleme her katlamada
sadece eğitime fit edilir).

\begin{lstlisting}
from sklearn.pipeline import Pipeline
from sklearn.preprocessing import StandardScaler
from sklearn.impute import SimpleImputer

pipeline = Pipeline([
    ('impute', SimpleImputer(strategy='median')),   # eksikleri doldur
    ('scale',  StandardScaler()),                    # standartlastir
    ('model',  LogisticRegression(max_iter=1000))    # tahmin et
])

pipeline.fit(X_train, y_train)    # tum zincir egitilir
pipeline.predict(X_test)          # tum zincir uygulanir
# GridSearchCV ile de dogrudan kullanilabilir
\end{lstlisting}

\ipucu{Gerçek projelerde \textbf{her zaman Pipeline kullan}. Ön işleme
adımlarını modelle birleştirmek; veri sızıntısını önler, kodu tekrardan kurtarır
ve modeli tek parça olarak kaydedip (\code{joblib.dump}) üretime almayı
kolaylaştırır. ``Manuel'' ön işleme, hataya en açık kısımdır.}

% =====================================================================
% #####################################################################
\gpart{6}{Uygulamalı Projeler: Gerçek Veri Setleri}
% #####################################################################

Buraya kadar araçları tek tek öğrendik. Şimdi hepsini \emph{gerçek, adı belli,
indirilebilir} veri setleri üzerinde birleştiriyoruz. Aşağıdaki projelerin her
biri kendi başına çalışır: veri seti bir tek satırla yüklenir (hepsi \code{seaborn}
veya \code{scikit-learn} ile birlikte gelir ya da internetten çekilir), böylece
kopyala--çalıştır--öğren yapabilirsin. Her projenin başında verinin kaynağı ve
resmî linki verilmiştir.

\bilgi{Kullanılan veri setleri ve kaynakları:
\begin{itemize}\setlength{\itemsep}{2pt}
\item \textbf{Palmer Penguins} --- \href{https://allisonhorst.github.io/palmerpenguins/}{allisonhorst.github.io/palmerpenguins} \\(seaborn: \code{sns.load\_dataset('penguins')})
\item \textbf{Titanic} --- \href{https://www.kaggle.com/c/titanic}{kaggle.com/c/titanic} \\(seaborn: \code{sns.load\_dataset('titanic')})
\item \textbf{Tips (Bahşiş)} --- \href{https://github.com/mwaskom/seaborn-data}{github.com/mwaskom/seaborn-data} \\(seaborn: \code{sns.load\_dataset('tips')})
\item \textbf{California Housing} --- \href{https://scikit-learn.org/stable/datasets/real_world.html\#california-housing-dataset}{scikit-learn --- California Housing} \\(\code{sklearn.datasets.fetch\_california\_housing})
\end{itemize}}

\ipucu{\code{sns.load\_dataset()} veriyi internetten çeker (bir kereliğine
indirir). İnternet yoksa aynı CSV'leri
\href{https://github.com/mwaskom/seaborn-data}{github.com/mwaskom/seaborn-data}
adresinden indirip \code{pd.read\_csv('penguins.csv')} ile açabilirsin.}

% =====================================================================
\gsec{Proje 1 --- Palmer Penguins: Keşifsel Veri Analizi}
% =====================================================================

\notk[Veri seti]{\textbf{Palmer Penguins}: Antarktika'daki üç penguen türünün
(Adelie, Chinstrap, Gentoo) gaga uzunluğu, kanat boyu, kütle gibi ölçümleri.
Iris'e modern bir alternatif. Kaynak:
\href{https://allisonhorst.github.io/palmerpenguins/}{allisonhorst.github.io/palmerpenguins}.}

Amaç: yeni bir veri setiyle ilk karşılaşmada izlenen \emph{keşif} adımlarını
(yükle $\to$ tanı $\to$ temizle $\to$ görselleştir $\to$ grupla) baştan sona
uygulamak.

\begin{lstlisting}
import seaborn as sns
import matplotlib.pyplot as plt
import pandas as pd

# 1. YUKLE
df = sns.load_dataset('penguins')

# 2. ILK BAKIS
print(df.shape)          # (344, 7)
print(df.head())
print(df.info())         # sutun tipleri + eksik degerler
print(df.describe())     # sayisal ozet
print(df['species'].value_counts())

# 3. EKSIK VERI ANALIZI VE TEMIZLEME
print(df.isna().sum())               # her sutunda kac eksik?
df = df.dropna().reset_index(drop=True)  # basit yol: eksikleri at
# (Alternatif: sayisallari medyan, kategorikleri mod ile doldurmak)
\end{lstlisting}

\begin{lstlisting}
# 4. GRUPLA & TOPLULASTIR: tur bazinda ortalama olcumler
ozet = df.groupby('species')[['bill_length_mm', 'flipper_length_mm',
                               'body_mass_g']].mean().round(1)
print(ozet)

# Tur + cinsiyet kiriliminda ortalama kutle
print(df.pivot_table(values='body_mass_g',
                     index='species', columns='sex', aggfunc='mean'))

# 5. GORSELLESTIR: dagilimlar ve iliskiler
sns.set_theme(style='whitegrid')
fig, axes = plt.subplots(1, 2, figsize=(12, 5))

sns.scatterplot(data=df, x='bill_length_mm', y='flipper_length_mm',
                hue='species', style='species', s=60, ax=axes[0])
axes[0].set_title('Gaga vs Kanat: turler ayrisir mi?')

sns.boxplot(data=df, x='species', y='body_mass_g', ax=axes[1])
axes[1].set_title('Ture gore vucut kutlesi')
fig.tight_layout(); plt.show()

# Tum sayisal ciftler + tur rengi (tek bakista kesif)
sns.pairplot(df, hue='species', height=1.8)
plt.show()
\end{lstlisting}

\ipucu{Dağılım grafiğinde türlerin net kümeler oluşturduğunu göreceksin ---
bu, bir sınıflandırma modelinin türü ölçümlerden \emph{kolayca} tahmin
edebileceğinin işaretidir. Keşif aşaması, modellemeden önce ``bu problem
çözülebilir mi?'' sorusunu yanıtlar.}

% =====================================================================
\gsec{Proje 2 --- Titanic: Temizleme ve Sınıflandırma}
% =====================================================================

\notk[Veri seti]{\textbf{Titanic}: 1912 kazasındaki yolcuların yaş, cinsiyet,
bilet sınıfı, ücret gibi bilgileri ve hayatta kalıp kalmadıkları. Makine
öğrenmesinin ``merhaba dünya''sı. Kaynak:
\href{https://www.kaggle.com/c/titanic}{kaggle.com/c/titanic}.}

Amaç: \emph{gerçekçi biçimde dağınık} bir veriyi (eksik yaşlar, kategorik
sütunlar) temizleyip özellik üretmek ve bir yolcunun hayatta kalıp kalmayacağını
tahmin eden bir model kurmak.

\begin{lstlisting}
import seaborn as sns
import pandas as pd

df = sns.load_dataset('titanic')
print(df[['survived','pclass','sex','age','fare','embarked']].head())
print(df.isna().sum())     # age ve deck cok eksik

# --- TEMIZLEME & OZELLIK MUHENDISLIGI ---
d = df[['survived','pclass','sex','age','sibsp','parch','fare','embarked']].copy()

# Eksik yasi, ayni pclass+sex grubunun MEDYANI ile doldur (akilli imputation)
d['age'] = d.groupby(['pclass','sex'])['age'].transform(
    lambda s: s.fillna(s.median()))
d['embarked'] = d['embarked'].fillna(d['embarked'].mode()[0])

# Yeni ozellikler: aile buyuklugu ve yalniz mi?
d['family_size'] = d['sibsp'] + d['parch'] + 1
d['is_alone'] = (d['family_size'] == 1).astype(int)

# Kategorikleri sayisala cevir (one-hot)
d = pd.get_dummies(d, columns=['sex','embarked'], drop_first=True)
print(d.head())
\end{lstlisting}

\begin{lstlisting}
# --- KESIF: hayatta kalma oranlari ---
print(df.groupby('sex')['survived'].mean())      # kadinlar cok daha yuksek
print(df.groupby('pclass')['survived'].mean())   # 1. sinif avantajli

import matplotlib.pyplot as plt
sns.barplot(data=df, x='pclass', y='survived', hue='sex')
plt.title('Sinif ve cinsiyete gore hayatta kalma orani'); plt.show()

# --- MODELLEME: Lojistik Regresyon ---
from sklearn.model_selection import train_test_split
from sklearn.linear_model import LogisticRegression
from sklearn.preprocessing import StandardScaler
from sklearn.pipeline import Pipeline
from sklearn.metrics import classification_report, confusion_matrix

X = d.drop(columns='survived')
y = d['survived']
X_tr, X_te, y_tr, y_te = train_test_split(
    X, y, test_size=0.2, random_state=42, stratify=y)

model = Pipeline([('scale', StandardScaler()),
                  ('clf', LogisticRegression(max_iter=1000))])
model.fit(X_tr, y_tr)

print(f"Test dogrulugu: {model.score(X_te, y_te):.3f}")
print(classification_report(y_te, model.predict(X_te)))
print(confusion_matrix(y_te, model.predict(X_te)))

# Hangi ozellik ne kadar etkili? (katsayilar)
katsayilar = pd.Series(model.named_steps['clf'].coef_[0], index=X.columns)
print(katsayilar.sort_values())
\end{lstlisting}

\uyari{Eksik \code{age} değerlerini \emph{tüm} sütunun ortalamasıyla doldurmak
kolaydır ama kabadır. Yukarıdaki gibi \code{groupby(['pclass','sex'])} medyanıyla
doldurmak, birinci sınıf bir kadının yaşını üçüncü sınıf bir erkeğinkinden farklı
tahmin ederek çok daha gerçekçi olur. İyi \emph{imputation} model başarısını
doğrudan artırır.}

% =====================================================================
\gsec{Proje 3 --- Tips: İstatistik ve Görselleştirme}
% =====================================================================

\notk[Veri seti]{\textbf{Tips (Bahşiş)}: Bir restoranda 244 hesabın tutarı,
bahşiş, gün, öğün, masa büyüklüğü. Görselleştirme ve ilişki analizi için ideal.
Kaynak:
\href{https://github.com/mwaskom/seaborn-data}{github.com/mwaskom/seaborn-data}.}

Amaç: sürekli iki değişken (hesap tutarı ve bahşiş) arasındaki ilişkiyi
görselleştirmek, korelasyonu ölçmek ve ``bahşiş yüzdesi neye bağlı?'' sorusunu
gruplama ile yanıtlamak.

\begin{lstlisting}
import seaborn as sns
import matplotlib.pyplot as plt

tips = sns.load_dataset('tips')
tips['tip_pct'] = tips['tip'] / tips['total_bill'] * 100   # bahsis yuzdesi

# --- KORELASYON ---
print(tips[['total_bill','tip','size','tip_pct']].corr().round(2))
# total_bill ile tip guclu pozitif; ama tip_pct ile zayif!

# --- GORSELLESTIR ---
sns.set_theme(style='whitegrid')
fig, ax = plt.subplots(1, 3, figsize=(15, 4))

# 1) Hesap vs bahsis + regresyon dogrusu, cinsiyete gore
sns.regplot(data=tips, x='total_bill', y='tip', ax=ax[0])
ax[0].set_title('Hesap arttikca bahsis artar')

# 2) Gune ve ogune gore ortalama bahsis yuzdesi
sns.barplot(data=tips, x='day', y='tip_pct', hue='time', ax=ax[1])
ax[1].set_title('Gun/ogune gore bahsis %')

# 3) Sigara icen/icmeyen dagilimi
sns.violinplot(data=tips, x='smoker', y='tip_pct', ax=ax[2])
ax[2].set_title('Sigara ve bahsis %')
fig.tight_layout(); plt.show()
\end{lstlisting}

\begin{lstlisting}
# --- GRUPLAMA ile ICGORU ---
print(tips.groupby('day')['tip_pct'].agg(['mean','count']).round(2))
print(tips.groupby(['time','sex'])['tip_pct'].mean().round(2))

# Basit hipotez: ogle mi aksam mi daha comert?
from scipy import stats
ogle = tips.loc[tips['time']=='Lunch', 'tip_pct']
aksam = tips.loc[tips['time']=='Dinner', 'tip_pct']
t, p = stats.ttest_ind(ogle, aksam, equal_var=False)
print(f"t={t:.2f}, p={p:.3f}  ->  p<0.05 ise fark anlamli")
\end{lstlisting}

\ipucu{Dikkat çekici bulgu: \code{total\_bill} ile \code{tip} güçlü ilişkilidir
(büyük hesap $\to$ büyük bahşiş), ama bahşiş \emph{yüzdesi} hesap tutarından
neredeyse bağımsızdır. Ham korelasyon ile ``oran'' bazlı analiz farklı hikâyeler
anlatır --- doğru soruyu sormak, doğru grafiği seçmek kadar önemlidir.}

% =====================================================================
\gsec{Proje 4 --- California Housing: Regresyon}
% =====================================================================

\notk[Veri seti]{\textbf{California Housing}: 1990 nüfus sayımından 20.640
bölgenin medyan gelir, ev yaşı, oda sayısı, konum bilgisi ve \emph{medyan ev
fiyatı}. Sürekli bir hedefi (fiyat) tahmin etmek için klasik regresyon verisi.
Kaynak:
\href{https://scikit-learn.org/stable/datasets/real_world.html\#california-housing-dataset}{scikit-learn --- California Housing}.}

Amaç: sürekli bir değeri (ev fiyatı) tahmin eden bir regresyon modeli kurmak,
doğru metriklerle (MAE, RMSE, $R^2$) değerlendirmek ve hangi özelliğin fiyatı en
çok belirlediğini bulmak.

\begin{lstlisting}
from sklearn.datasets import fetch_california_housing
import pandas as pd

data = fetch_california_housing(as_frame=True)
df = data.frame
print(df.head())
print(df.describe())
# Hedef: MedHouseVal (100.000 dolar biriminde medyan ev degeri)

# Hangi ozellik fiyatla en iliskili?
print(df.corr()['MedHouseVal'].sort_values(ascending=False))
# MedInc (medyan gelir) acik ara en guclu belirleyici
\end{lstlisting}

\begin{lstlisting}
from sklearn.model_selection import train_test_split, cross_val_score
from sklearn.preprocessing import StandardScaler
from sklearn.linear_model import LinearRegression
from sklearn.ensemble import RandomForestRegressor
from sklearn.pipeline import Pipeline
from sklearn.metrics import mean_absolute_error, mean_squared_error, r2_score
import numpy as np

X = df.drop(columns='MedHouseVal')
y = df['MedHouseVal']
X_tr, X_te, y_tr, y_te = train_test_split(X, y, test_size=0.2, random_state=42)

# Iki modeli karsilastir: dogrusal (basit) vs random forest (guclu)
modeller = {
    'Linear': Pipeline([('sc', StandardScaler()), ('m', LinearRegression())]),
    'RandomForest': RandomForestRegressor(n_estimators=100, random_state=42,
                                          n_jobs=-1),
}
for ad, model in modeller.items():
    model.fit(X_tr, y_tr)
    tahmin = model.predict(X_te)
    mae = mean_absolute_error(y_te, tahmin)
    rmse = np.sqrt(mean_squared_error(y_te, tahmin))
    r2 = r2_score(y_te, tahmin)
    print(f"{ad:14s} MAE={mae:.3f}  RMSE={rmse:.3f}  R2={r2:.3f}")

# En iyi modelin ozellik onemleri
rf = modeller['RandomForest']
onem = pd.Series(rf.feature_importances_, index=X.columns)
print(onem.sort_values(ascending=False))
\end{lstlisting}

\ipucu{İki modeli karşılaştırınca Random Forest'ın doğrusal modeli belirgin
biçimde geçtiğini göreceksin --- çünkü fiyat, özelliklerle doğrusal olmayan ve
etkileşimli bir ilişki içinde (örn. konumun etkisi gelire bağlı değişir).
Regresyonda da ``basit bir temel çizgi (Linear) kur, sonra güçlü bir modelle
kıyasla'' altın kuraldır.}

% =====================================================================
\gsec{Kapanış: Uçtan Uca Bir Proje}
% =====================================================================

Öğrendiğimiz her şeyi birleştiren tipik bir veri bilimi projesinin iskeleti
şöyledir. Gerçek beceri, bu adımları veriye göre uyarlamaktır.

\begin{lstlisting}
import pandas as pd
from sklearn.model_selection import train_test_split, cross_val_score
from sklearn.pipeline import Pipeline
from sklearn.preprocessing import StandardScaler
from sklearn.ensemble import RandomForestClassifier
from sklearn.metrics import classification_report

# 1. VERIYI YUKLE
df = pd.read_csv('veri.csv')

# 2. KESIF & TEMIZLE
df = df.drop_duplicates()
df['age'] = df['age'].fillna(df['age'].median())
print(df.describe())

# 3. OZELLIK & HEDEF AYIR
X = df.drop(columns=['target'])
y = df['target']

# 4. EGITIM/TEST AYIR
X_train, X_test, y_train, y_test = train_test_split(
    X, y, test_size=0.2, random_state=42, stratify=y)

# 5. PIPELINE KUR & EGIT
pipeline = Pipeline([
    ('scale', StandardScaler()),
    ('model', RandomForestClassifier(n_estimators=100, random_state=42))
])
pipeline.fit(X_train, y_train)

# 6. DEGERLENDIR
scores = cross_val_score(pipeline, X_train, y_train, cv=5)
print(f"CV dogruluk: {scores.mean():.3f}")
print(classification_report(y_test, pipeline.predict(X_test)))

# 7. KAYDET
import joblib
joblib.dump(pipeline, 'model.pkl')
\end{lstlisting}

\gsub{Özet: Doğru Aracı Seçmek}

\begin{center}
\begin{tabularx}{\textwidth}{@{}>{\bfseries\color{anaMavi}}l X@{}}
\toprule
\rowcolor{acikMavi}\color{anaMavi}\textbf{Görev} & \color{anaMavi}\textbf{Araç / Yöntem}\\
\midrule
Sayısal hesaplama, diziler & NumPy\\
\rowcolor{acikGri} Tablo verisi, temizleme & Pandas\\
Hızlı grafik & Pandas \code{.plot()}, Matplotlib\\
\rowcolor{acikGri} İstatistiksel grafik & Seaborn\\
Sürekli değer tahmini & Regresyon (Linear, RandomForest)\\
\rowcolor{acikGri} Kategori tahmini & Sınıflandırma (Logistic, RF, SVM)\\
Grup bulma (etiketsiz) & K-Means, boyut indirgeme (PCA)\\
\rowcolor{acikGri} Güvenilir değerlendirme & Çapraz doğrulama, Pipeline\\
\bottomrule
\end{tabularx}
\end{center}

\ipucu{Veri biliminde başarının anahtarı fancy modeller değil, \emph{temiz veri
ve doğru değerlendirmedir}. Zamanının çoğu adım 2--3'te (temizleme, keşif)
geçecek --- ve bu normaldir. Basit bir modelle başla, temel bir başarı çizgisi
(baseline) kur, sonra iyileştir. Bol gerçek veri setiyle pratik yapmak, bu
sezgiyi geliştiren tek yoldur.}

\vfill
\begin{center}
{\color{ortaGri}\rule{0.5\linewidth}{0.6pt}}\\[6pt]
{\color{ortaGri}\small Rehberin sonu --- İyi analizler!}
\end{center}

\end{document}
